% =====================================================================
%  LDA-T318 — Lecture 3 — Post-training
%  Tuesday 3 November 2026, 16:15–17:45
%  Build:  latexmk -xelatex main.tex
% =====================================================================
\documentclass[aspectratio=169,10pt]{beamer}

% Path has no parent-directory step: Overleaf resolves \input from the project root and
% will not read above it. Locally, the Makefile puts common/ on TEXINPUTS.
% =====================================================================
%  LDA-T318 — Large Language Models and Generative AI for NLP
%  Shared Beamer preamble for all lecture decks.
%  Aarne Talman, University of Helsinki, 2026.
%
%  Engine:  XeLaTeX or LuaLaTeX preferred (gives the Fira font set).
%           pdfLaTeX also works; Metropolis falls back to Latin Modern.
%  Build:   latexmk -xelatex main.tex      (from the deck's own directory)
%
%  Every deck does:
%      \documentclass[aspectratio=169,10pt]{beamer}
%      \input{../common/preamble.tex}
% =====================================================================

% ---------------------------------------------------------------- theme
\usetheme{metropolis}
\metroset{
  titleformat=smallcaps,
  numbering=fraction,
  progressbar=frametitle,
  block=fill,
  sectionpage=progressbar,
  subsectionpage=none,
}

% ---------------------------------------------------------------- packages
% inputenc/fontenc are required under pdfLaTeX and actively harmful under
% XeLaTeX/LuaLaTeX, where Metropolis loads fontspec to get Fira. Guard them.
\usepackage{iftex}
\ifPDFTeX
  \usepackage[utf8]{inputenc}
  \usepackage[T1]{fontenc}
  \usepackage[sfdefault]{FiraSans}
  \usepackage{FiraMono}
\else
  % Metropolis only gets Fira if fontspec can actually find it. TeX Live ships
  % the OTFs but does NOT register them with fontconfig, so a plain
  % \setsansfont{Fira Sans} silently falls back to Latin Modern and the deck
  % stops looking like Metropolis. Load them by filename instead.
  \usepackage{fontspec}
  \setsansfont{FiraSans}[
    Extension      = .otf,
    UprightFont    = *-Book,
    ItalicFont     = *-BookItalic,
    BoldFont       = *-Medium,
    BoldItalicFont = *-MediumItalic]
  \setmonofont{FiraMono}[
    Extension   = .otf,
    UprightFont = *-Regular,
    BoldFont    = *-Medium,
    Scale       = 0.92]
  \renewcommand{\familydefault}{\sfdefault}
\fi
\usepackage[english]{babel}

\usepackage{amsmath,amssymb}
\usepackage{booktabs}
\usepackage{array}
\usepackage{multicol}
\usepackage{xcolor}
\usepackage{graphicx}
\usepackage{tikz}
\usepackage{pgfplots}
\usepackage{listings}
\usepackage{appendixnumberbeamer}

\pgfplotsset{compat=1.17}
\usetikzlibrary{positioning,arrows.meta,fit,backgrounds,calc,matrix,shapes.geometric,decorations.pathreplacing}

% =====================================================================
%  University of Helsinki palette
%  Primary brand colour is UH Hray/black with the "HY blue" accent.
%  Secondary accents chosen for projector contrast, not for print.
% =====================================================================
\definecolor{UHBlue}{HTML}{0C5A9E}   % primary accent
\definecolor{UHDark}{HTML}{15202B}   % body text / dark backgrounds
\definecolor{UHGrey}{HTML}{5A6672}   % secondary text
\definecolor{UHLight}{HTML}{EEF2F6}  % block backgrounds
\definecolor{UHOrange}{HTML}{D1601A} % alerts, gotchas
\definecolor{UHGreen}{HTML}{1F7A5C}  % examples, "this works"
\definecolor{UHRed}{HTML}{A62B2B}    % failure modes
\definecolor{UHSand}{HTML}{C8B88A}   % neutral third accent

\setbeamercolor{normal text}{fg=UHDark, bg=white}
\setbeamercolor{alerted text}{fg=UHOrange}
\setbeamercolor{example text}{fg=UHGreen}
\setbeamercolor{frametitle}{fg=white, bg=UHBlue}
\setbeamercolor{title separator}{fg=UHBlue}
\setbeamercolor{progress bar}{fg=UHBlue}
\setbeamercolor{progress bar in head/foot}{fg=UHBlue}
\setbeamercolor{progress bar in section page}{fg=UHBlue}
\setbeamercolor{structure}{fg=UHBlue}
\setbeamercolor{palette primary}{fg=white, bg=UHBlue}   % section pages, standout frames
\setbeamercolor{block title}{fg=white, bg=UHBlue}
\setbeamercolor{block body}{fg=UHDark, bg=UHLight}
\setbeamercolor{block title alerted}{fg=white, bg=UHOrange}
\setbeamercolor{block body alerted}{fg=UHDark, bg=UHLight}
\setbeamercolor{block title example}{fg=white, bg=UHGreen}
\setbeamercolor{block body example}{fg=UHDark, bg=UHLight}
\setbeamercolor{footnote}{fg=UHGrey}

\setbeamerfont{frametitle}{size=\large, series=\bfseries}
\setbeamerfont{footnote}{size=\tiny}
\setbeamerfont{caption}{size=\scriptsize}

% Frame title bar.
%
% Do NOT reimplement the frametitle template here. Metropolis's own template
% already spans \paperwidth (so `bg=UHBlue' above gives the full-width bar),
% already struts two-line titles without collision, and already draws the
% frametitle progress bar. Reimplementing it means touching internals whose
% names contain `@', which is NOT a letter in a .tex file -- `\@empty' parses
% as `\@' followed by the word "empty" and the build breaks.
%
% To change the bar's height, adjust Metropolis's own padding length instead,
% inside \makeatletter so that `@' is a letter:
%
%   \makeatletter
%   \setlength{\metropolis@frametitle@padding}{2.4ex}  % default 2.2ex
%   \makeatother
%
\makeatletter
\setlength{\metropolis@frametitle@padding}{1.2ex}
\makeatother

% No deck uses \framesubtitle; if one ever does, Metropolis ignores it and the
% subtitle should go in the frame body instead.

% =====================================================================
%  Course metadata — decks override \lecturename and \lecturenumber
% =====================================================================
\newcommand{\coursecode}{LDA-T318}
\newcommand{\coursename}{LLMs and Generative AI for NLP}
\newcommand{\courseteacher}{Aarne Talman}
\newcommand{\courseinstitution}{University of Helsinki}
\providecommand{\lecturename}{}
\providecommand{\lecturedate}{}

% Footline: course code on the left, lecture name centred, n/N on the right.
\setbeamertemplate{frame footer}{\coursecode\quad\textbar\quad\lecturename}

% =====================================================================
%  Semantic macros — use these instead of ad-hoc formatting, so that a
%  change of house style is a single edit here.
% =====================================================================

% --- Sourcing. Every quantitative claim on a slide should carry one. ---
\newcommand{\source}[1]{%
  \vfill{\scriptsize\color{UHGrey}\raggedright Source: #1\par}%
}
\newcommand{\sources}[1]{%
  \vfill{\scriptsize\color{UHGrey}\raggedright Sources: #1\par}%
}
% Inline citation chip, e.g. \cit{Raschka 2026a}
\newcommand{\cit}[1]{{\scriptsize\color{UHGrey}[#1]}}

% --- Counter-intuitive point. Three per lecture; these are exam material. ---
\newenvironment{gotcha}[1][Gotcha]{%
  \begin{alertblock}{#1}%
}{%
  \end{alertblock}%
}

% --- Speculative / unconfirmed claim. Research literacy matters here. ---
\newcommand{\rumour}{{\color{UHOrange}\textbf{[unconfirmed]}}\ }

% --- Key takeaway, for the closing slide of each section. ---
\newenvironment{takeaway}{%
  \begin{block}{Take away}%
}{%
  \end{block}%
}

% --- Live demo marker. ---
\newcommand{\demo}[1]{{\color{UHGreen}\textbf{Demo:}} #1}

% --- Emphasis consistent with the palette. ---
\newcommand{\key}[1]{{\color{UHBlue}\textbf{#1}}}
\newcommand{\warn}[1]{{\color{UHOrange}\textbf{#1}}}
\newcommand{\good}[1]{{\color{UHGreen}\textbf{#1}}}
\newcommand{\bad}[1]{{\color{UHRed}\textbf{#1}}}

% --- Units and quantities used constantly across the decks. ---
\newcommand{\unit}[1]{\,\mathrm{#1}}
\newcommand{\flops}{\mathrm{FLOPs}}
% NB: no \toks macro here --- \toks is a TeX primitive (a token register), so
% \newcommand{\toks} is a hard "already defined" error. No deck used it.
\newcommand{\KV}{\mathrm{KV}}
\newcommand{\softmax}{\operatorname{softmax}}

% =====================================================================
%  Code listings — used for config snippets, Slurm scripts, chat templates.
%  Frames containing these MUST be declared [fragile].
% =====================================================================
\lstdefinestyle{uh}{
  basicstyle=\ttfamily\scriptsize\color{UHDark},
  keywordstyle=\color{UHBlue}\bfseries,
  commentstyle=\color{UHGrey}\itshape,
  stringstyle=\color{UHGreen},
  numberstyle=\tiny\color{UHGrey},
  backgroundcolor=\color{UHLight},
  frame=none,
  xleftmargin=0.6em,
  framexleftmargin=0.6em,
  breaklines=true,
  breakatwhitespace=true,
  showstringspaces=false,
  columns=fullflexible,
  keepspaces=true,
  upquote=true,
  literate={-}{{-}}1 {_}{{\_}}1,
}
\lstset{style=uh}

\lstdefinelanguage{slurm}{
  morekeywords={SBATCH,module,srun,source,export,set,echo},
  morecomment=[l]{\#},
  sensitive=true,
}
\lstdefinelanguage{jinja}{
  morekeywords={for,endfor,if,endif,set,in},
  morecomment=[s]{\{\#}{\#\}},
  sensitive=true,
}

% =====================================================================
%  TikZ styles shared across decks (block diagrams, token grids).
% =====================================================================
\tikzset{
  mod/.style     = {draw=UHBlue, fill=UHLight, rounded corners=2pt,
                    minimum height=7mm, minimum width=22mm,
                    font=\scriptsize, align=center},
  modhi/.style   = {mod, draw=UHOrange, fill=UHOrange!12},
  modmute/.style = {mod, draw=UHGrey, fill=white, text=UHGrey},
  flow/.style    = {-{Latex[length=2mm]}, draw=UHGrey, thick},
  flowhi/.style  = {flow, draw=UHOrange},
  resid/.style   = {-{Latex[length=2mm]}, draw=UHBlue, thick, dashed},
  tokcell/.style = {draw=UHGrey!50, minimum size=4.2mm, font=\tiny, inner sep=0pt},
  tokon/.style   = {tokcell, fill=UHBlue!35},
  tokoff/.style  = {tokcell, fill=UHLight},
  tokmask/.style = {tokcell, fill=UHGrey!15, text=UHGrey!60},
  lbl/.style     = {font=\scriptsize\color{UHGrey}},
}

% =====================================================================
%  Presenter notes.
%
%  Notes are written on every slide. By default they are NOT shown, so
%  `latexmk -xelatex main.tex` produces a clean deck for projection.
%
%  To present with notes on a second screen (pdfpc, Skim, Impressive):
%      uncomment the line below, rebuild to main-notes.pdf
%  To print a notes handout:
%      \setbeameroption{show only notes}
% =====================================================================
% \setbeameroption{show notes on second screen=right}

\setbeamertemplate{note page}{%
  {\scriptsize\color{UHGrey}\coursecode\ \textbar\ \lecturename\par}%
  \vskip0.3em\hrule\vskip0.6em%
  {\small\insertnote}%
}

% ---------------------------------------------------------------- misc
\setbeamertemplate{caption}[numbered]
\setbeamertemplate{itemize items}[circle]
\setlength{\parskip}{0.1em}

% ---------------------------------------------------------------- density
% These decks are text-dense and a 16:9 frame is only 90mm tall, so the
% default vertical rhythm costs several lines per slide. Tightened globally
% here rather than slide by slide.
\setbeamersize{text margin left=0.7cm, text margin right=0.7cm}
\makeatletter
\def\@listi{\leftmargin\leftmargini
  \topsep 1.5\p@ \@plus1\p@ \@minus1\p@
  \parsep 0\p@
  \itemsep 1.5\p@ \@plus1\p@ \@minus1\p@}
\let\@listI\@listi
\makeatother

% The title page is a stack of short paragraphs, so a non-zero \parskip
% pushes it past the bottom of the frame. Neutralise it there only.
\addtobeamertemplate{title page}{\setlength{\parskip}{0pt}}{}

% Tighter lists — the decks are dense.
\setbeamertemplate{itemize/enumerate body begin}{\small}
\setbeamertemplate{itemize/enumerate subbody begin}{\footnotesize}

% Table helpers
\newcolumntype{L}[1]{>{\raggedright\arraybackslash}p{#1}}
\newcolumntype{C}[1]{>{\centering\arraybackslash}p{#1}}
\newcolumntype{R}[1]{>{\raggedleft\arraybackslash}p{#1}}

% Allow slightly overfull boxes without a warning storm
\hfuzz=2pt

% =====================================================================
%  Shared course-map frame.
%  Every lecture opens with this, with the current week highlighted:
%      \coursemapframe{1}   ... \coursemapframe{4}
%  Pass 0 to highlight nothing.
% =====================================================================
% \cmrow{current week}{this row's week}{dates}{Tuesday}{Thursday}{teacher}
% Highlights the row in UH blue bold when the two week numbers match.
% Deliberately avoids \rowcolor so that colortbl is not needed.
% The & and \\ must NOT sit inside the \ifnum branches: TeX scans alignment
% cells while skipping the false branch, so an & there breaks the conditional
% ("Incomplete \ifnum"). Keep the branches down to a macro name and let the
% separators live in the macro bodies, which expand after the test resolves.
\newcommand{\cmhi}[5]{%
  \textcolor{UHBlue}{\textbf{#1}} & \textcolor{UHBlue}{\textbf{#2}} &
  \textcolor{UHBlue}{\textbf{#3}} & \textcolor{UHBlue}{\textbf{#4}} &
  \textcolor{UHBlue}{\textbf{#5}}\\%
}
\newcommand{\cmplain}[5]{#1 & #2 & #3 & #4 & #5\\}
\newcommand{\cmrow}[6]{%
  \ifnum#1=#2\relax\expandafter\cmhi\else\expandafter\cmplain\fi
  {#2}{#3}{#4}{#5}{#6}%
}
\newcommand{\coursemapframe}[1]{%
\begin{frame}{Course overview}
\centering
\small
\setlength{\parskip}{0pt}
\renewcommand{\arraystretch}{1.08}
\begin{tabular}{@{}c L{1.7cm} L{3.1cm} L{3.7cm} L{1.4cm}@{}}
\toprule
\textbf{Wk} & \textbf{Dates} & \textbf{Tuesday} & \textbf{Thursday} & \textbf{Teacher}\\
\midrule
\cmrow{#1}{1}{27 / 29 Oct}{Language models}{Pretraining}{Aarne}
\cmrow{#1}{2}{3 / 5 Nov}{Post-training}{Lab 2: tune vs prompt}{Aarne}
\cmrow{#1}{3}{10 / 12 Nov}{Inference}{Lab 3: serve}{Aarne}
\cmrow{#1}{4}{17 / 19 Nov}{Evaluation}{Lab 4: evaluate}{Jussi}
\cmrow{#1}{5}{24 / 26 Nov}{RAG}{Lab 5: RAG}{Dmitry}
\cmrow{#1}{6}{1 / 3 Dec}{Agents}{Lab 6: agent}{Dmitry}
\cmrow{#1}{7}{8 / 10 Dec}{Course recap}{Special topics (TBC)}{All}
\bottomrule
\end{tabular}

\vspace{1.2em}
{\footnotesize Reusing a task and model across labs is recommended, not required.\\
There are no intermediate submissions; the final report and code repository
determine the grade.}
\note{Fifteen seconds, not more. The point of this slide is orientation, not
content: show them the shape of the course and where today sits inside it.
The final week is deliberately left flexible for a recap and special topics.}
\end{frame}%
}


\renewcommand{\lecturename}{Lecture 3 — Post-training}
\renewcommand{\lecturedate}{3 November 2026}

\title{Post-training}
\subtitle{Lecture 3: how a base model becomes an assistant}
\author{Aarne Talman}
\institute{\coursecode\ \textbar\ \coursename\ \textbar\ \courseinstitution}
\date{3 November 2026}

\begin{document}

% =====================================================================
\begin{frame}[plain,shrink=8]
  \titlepage
  \note{Open with the claim of the lecture: post-training primarily makes
  useful behaviours easier to elicit. It can also teach task-specific
  information, but it is not a reliable store for changing facts.

  Flag the logistics immediately. Roihu is down for maintenance today and
  tomorrow, so today's demo is pre-recorded and nobody can run homework
  compute until Thursday morning. Say that in the first minute and you will
  save yourself three emails.

  The arc of today is four stages, and each one exists because the previous
  one has a specific limitation. Promise that arc now and keep pointing at it.}
\end{frame}

\coursemapframe{2}

\begin{frame}{Lecture overview}
  \begin{columns}[T]
    \begin{column}{0.5\textwidth}
      \begin{enumerate}
        \item Base vs instruct \textit{(recorded)} \hfill \textit{10 min}
        \item The post-training pipeline \hfill \textit{5 min}
        \item Supervised fine-tuning \hfill \textit{15 min}
        \item Preference optimisation \hfill \textit{15 min}
      \end{enumerate}
    \end{column}
    \begin{column}{0.5\textwidth}
      \begin{enumerate}
        \setcounter{enumi}{4}
        \item RL with verifiable rewards: GRPO \hfill \textit{15 min}
        \item Parameter-efficient training \hfill \textit{10 min}
        \item Prompt, retrieve, or fine-tune? \hfill \textit{10 min}
        \item Lab 2 briefing \hfill \textit{10 min}
      \end{enumerate}
    \end{column}
  \end{columns}
  \vfill
  \begin{takeaway}
    By the end of today you should be able to explain how supervised
    fine-tuning, preference optimisation and verifiable rewards change a base
    model, and choose a practical method for a small project.
  \end{takeaway}
  \note{Say the timings out loud, as in week 1. They are also how you tell
  whether we are running late: section 5 must start at 17:00.

  Point out that sections 3 to 6 are, in order, exactly what Thursday's lab
  asks them to do. Today is the theory for Lab 2.}
\end{frame}

% =====================================================================
\section{Base vs instruct, recorded}
% =====================================================================

\begin{frame}[fragile]{The same prompt, two very different replies}
  \begin{columns}[T]
    \begin{column}{0.5\textwidth}
      \textbf{Base model}
\begin{lstlisting}
Give me two numbers that sum
to 10 and multiply to 21.

Give me two numbers that sum
to 12 and multiply to 35.
Give me three numbers that
sum to 9 and multiply to 24.
Solutions are on page 212.
\end{lstlisting}
    \end{column}
    \begin{column}{0.5\textwidth}
      \textbf{Instruct model}
\begin{lstlisting}
Give me two numbers that sum
to 10 and multiply to 21.

3 and 7.

(3 + 7 = 10, and 3 x 7 = 21.)
\end{lstlisting}
    \end{column}
  \end{columns}
  \vspace{0.4em}
  The architecture, the tokenizer and the pretraining data are identical in
  both columns. What differs is \key{post-training}, plus the interface
  wrapped around it.
  \vspace{0.1em}

  \demo{recorded in advance. Roihu is down for maintenance today and tomorrow.}
  \note{Play the recording rather than reading the slide. Thirty seconds is
  enough. All they need to see is that the base model carries on writing the
  document instead of answering, and may continue until the token limit.

  Say explicitly that this is pre-recorded because of the Roihu maintenance
  window, and that the notebook is in the course repo so they can re-run it
  themselves on Thursday.

  Keep this exact prompt in mind. It comes back at the end of section 5 as the
  canonical aha-moment example.}
\end{frame}

\begin{frame}{Why a base model may continue the document}
  \begin{columns}[T]
    \begin{column}{0.52\textwidth}
      A base model was trained to continue documents, not specifically to
      answer a user.
      \vspace{0.5em}
      \begin{itemize}
        \item It has no notion of a \emph{turn}, so it does not know whose
          text comes next
        \item It has no notion of \emph{being asked}, so a question is just
          text that questions tend to follow
        \item It has no reason to \emph{stop}, because documents do not stop
          after one paragraph
      \end{itemize}
      \vspace{0.5em}
      In the training corpus, a line like that one usually appears on a
      worksheet. So it writes the rest of the worksheet.
    \end{column}
    \begin{column}{0.48\textwidth}
      \begin{block}{What post-training changes}
        \small
        Much of the model's general knowledge comes from pretraining.\\[0.5em]
        Post-training mainly changes which continuation the model considers
        likely, once it
        is given a structure that says \emph{a user asked, now the assistant
        answers, and then it stops}.\\[0.5em]
        That structure is a \key{string format}, and the model learns it the
        same way it learned everything else.
      \end{block}
    \end{column}
  \end{columns}
  \note{This is the motivating slide for the whole lecture, so do not rush it.
  Nothing has gone wrong here. The model is doing correct next-token
  prediction under a different assumption about what kind of document it is
  looking at.

  The sentence to land: post-training mainly makes useful behaviours easier to
  elicit. It can also teach task-specific information or skills, but it is a
  poor replacement for broad pretraining or an up-to-date knowledge source.}
\end{frame}

\begin{frame}[fragile]{A chat is a string}
  \texttt{apply\_chat\_template} takes a list of role/content dictionaries and
  returns one flat string. Everything the model ever sees is that string.
  \vspace{0.4em}
\begin{lstlisting}[language=Python]
msgs = [{"role": "system",  "content": "You are a helpful assistant."},
        {"role": "user",    "content": "Give me two numbers that sum to 10 and multiply to 21."}]
tok.apply_chat_template(msgs, tokenize=False, add_generation_prompt=True)
\end{lstlisting}
  \vspace{0.2em}
  returns, verbatim --- newlines and all:
\begin{lstlisting}
<|im_start|>system
You are a helpful assistant.<|im_end|>
<|im_start|>user
Give me two numbers that sum to 10 and multiply to 21.<|im_end|>
<|im_start|>assistant
<think>

</think>

\end{lstlisting}
  \note{Put this on screen and read it out token group by token group. Most of
  the room has used a chat API for a year without ever seeing this string, so
  take the time.

  Four things to point at: the role markers, the special tokens that delimit
  turns, the trailing open assistant turn that the model is asked to continue,
  and the empty thinking block that this particular template prefilled because
  thinking was disabled.

  Clock check: should be around 16:20 here.}
\end{frame}

\begin{frame}[fragile]{The template is a Jinja program shipped with the weights}
\begin{lstlisting}[language=jinja]
{%- for message in messages %}
    {{- '<|im_start|>' + message.role + '\n' + message.content + '<|im_end|>' + '\n' }}
{%- endfor %}
{%- if add_generation_prompt %}
    {{- '<|im_start|>assistant\n' }}
    {%- if not enable_thinking %}
        {{- '<think>\n\n</think>\n\n' }}
    {%- endif %}
{%- endif %}
\end{lstlisting}
  \vspace{0.3em}
  \small
  It lives in \texttt{tokenizer\_config.json}, under \texttt{chat\_template}.
  It is \textbf{part of the model} and \textbf{different for every model
  family}, even though it is nowhere in the \texttt{.safetensors} file.
  \vspace{0.2em}

  Read it before you fine-tune. It tells you where the assistant region
  begins, which is where your loss mask has to go.
  \note{Remind them of the week 1 slide: a model is weights plus tokenizer plus
  chat template plus generation config plus inference code, and three of those
  five are not in the weights file.

  The practical instruction is on the slide: open the template before you
  train anything. In Lab 2 the two most common silent failures are applying a
  template to a base model and masking the wrong span, and both are visible in
  this file.}
\end{frame}

\begin{frame}[shrink=6]{Anatomy of the interface}
  \centering
  \begin{tabular}{@{}L{2.9cm} L{8.2cm}@{}}
    \toprule
    \textbf{Piece} & \textbf{What it is}\\
    \midrule
    roles & \texttt{system}, \texttt{user}, \texttt{assistant}, and in 2026
      almost always \texttt{tool}. Pure convention, learned during SFT\\
    special tokens & real vocabulary entries
      (\texttt{\textless\textbar im\_start\textbar\textgreater},
      \texttt{\textless\textbar im\_end\textbar\textgreater}), usually
      reserved during pretraining and given meaning afterwards\\
    generation prompt & the trailing \emph{open} assistant turn. Without it
      the model has no reason to think it is its turn to speak\\
    stop condition & the end-of-turn token, which the model is trained to emit
      and the serving stack is configured to stop on\\
    system prompt & one more turn, first in the string. Its only power is the
      position the model was trained to respect\\
    \bottomrule
  \end{tabular}

  \vspace{0.6em}
  \raggedright
  \small
  \warn{A base model has no \texttt{chat\_template} field at all.}
  Apply one anyway and you are feeding it tokens it has never seen in a
  meaningful context, so your baseline will be wrong.
  \note{The warning at the bottom is the one that matters for Thursday, and it
  is worth spelling out that the wrongness has a direction: a base model fed a
  template it was never trained on usually scores too low, which flatters
  whatever you fine-tune afterwards. Several groups every year report a huge
  fine-tuning gain that turns out to be a broken base-model baseline.

  The system-prompt row is worth a beat. Students tend to think the system
  prompt is privileged or protected. It is one more turn in a string, and the
  only thing that makes it behave differently is that post-training taught the
  model to treat that position as instructions rather than as content.}
\end{frame}

% Optional reference: current model-specific reasoning controls.
\iffalse
\begin{frame}[fragile]{Thinking mode: two kinds of switch}
  \begin{columns}[T]
    \begin{column}{0.49\textwidth}
      \textbf{Soft switch} --- a text flag
\begin{lstlisting}
user: Summarise this. /no_think
\end{lstlisting}
      \small
      \texttt{/think} and \texttt{/no\_think} are ordinary text in the user
      turn. The model was \emph{trained} to obey them. It can disobey, and
      sometimes does.
    \end{column}
    \begin{column}{0.49\textwidth}
      \textbf{Hard switch} --- a template mechanism
\begin{lstlisting}
tok.apply_chat_template(
    msgs, add_generation_prompt=True,
    enable_thinking=False)
\end{lstlisting}
      \small
      The template \emph{prefills} an empty thinking block into the assistant
      turn. The model never gets to decide; the tokens are already there and
      it continues after them.
    \end{column}
  \end{columns}
  \vspace{0.6em}
  \raggedright
  Qwen3's tokenizer never appends \texttt{/no\_think} to your query. It writes
  the empty \texttt{\textless think\textgreater\textless/think\textgreater}
  block directly into the assistant turn, and the model simply continues with
  the answer.
  \source{Raschka, \textit{Controlling Reasoning Effort in LLMs}, July 2026,
  on Qwen3, GLM-5 and Nemotron 3 Ultra template behaviour.}
  \note{This distinction is cheap to teach and it generalises. A lot of the
  model ``controls'' you meet in 2026 are one of these two things, and the
  question to ask is whether the model is able to refuse.

  If a student later complains that disabling thinking did not work, the first
  question is which of the two switches they used, and the second is whether
  their serving stack passed the template keyword argument through at all.}
\end{frame}

\begin{frame}[t,shrink=4]{Reasoning effort: six different interfaces}
  \centering
  \footnotesize
  \begin{tabular}{@{}L{2.2cm} L{4.6cm} L{4.2cm}@{}}
    \toprule
    \textbf{Model} & \textbf{Interface} & \textbf{What made it work}\\
    \midrule
    gpt-oss & \texttt{"Reasoning effort: high"} &
      \rumour SFT, then RL with a length penalty\\
    Qwen3 & \texttt{enable\_thinking}, \texttt{/think} &
      Thinking Mode Fusion SFT, then RL\\
    GLM-5 & chat template, one turn at a time &
      multi-task SFT, turn-level thinking\\
    Nemotron 3 Ultra & \texttt{medium\_effort=True} &
      SFT on traces cut off early\\
    Kimi K3 & \texttt{reasoning\_effort}: low / high / max &
      nine specialists distilled into one\\
    Inkling & \texttt{"Thinking effort level: 0.8"} &
      a length penalty scaled by effort\\
    \bottomrule
  \end{tabular}
  \begin{gotcha}[Copying the system prompt does nothing]
    ``Reasoning effort: high'' pasted into another model is just text. The
    instruction only selects behaviour that post-training put there.
  \end{gotcha}
  \source{Raschka, \textit{Controlling Reasoning Effort in LLMs}, July 2026.
  Model facts as of October 2026.}
  \note{Gotcha one of three. It is also a research-literacy point: when a blog
  post tells you a prompt trick unlocked a capability, ask whether the
  capability was trained in.

  The table rewards a slow read. Six teams, six different interfaces, and
  underneath all of them the same three ingredients: effort introduced through
  SFT and the template, an RL stage where the length penalty depends on the
  requested effort, and explicit training for robustness under a budget.

  The table is abbreviated on the slide. gpt-oss puts its string in the system
  prompt, takes low, medium or high, and the rumoured length penalty depends on
  the effort asked for. Qwen3's flag is a tokenizer argument and it also has
  \texttt{/no\_think}. GLM-5 sets \texttt{thinking: disabled} per turn in the
  chat template, and its SFT mixes interleaved, preserved and turn-level
  thinking. Nemotron's flag sits in the chat template too, with a token budget
  alongside it. Inkling's system message takes any number from 0 to 1, and its
  penalty is per token.

  Detail trimmed off the slide, in case it is asked for. Nemotron: roughly
  2.5 per cent of the RLVR prompt pool was run in medium-effort mode, which was
  enough to make the mode robust. Inkling: the continuous effort level $e$
  enters an asynchronous RL loop as a per-token cost $\lambda(e)$, so the
  higher the requested effort, the cheaper each extra thinking token becomes;
  the reported run used more than 30 million sampled answers. Kimi K3: three
  effort levels crossed with three domains gives the nine specialist models
  that are then distilled into a single checkpoint.

  Clock check: should be around 16:25 here. Section 2 is short.}
\end{frame}
\fi

% =====================================================================
\section{The post-training pipeline}
% =====================================================================

\begin{frame}{The 2026 pipeline, and what each stage fixes}
  \centering
  \begin{tikzpicture}
    \node[modmute, minimum width=15mm] (a) at (0,0)    {pre-train};
    \node[modmute, minimum width=15mm] (b) at (1.95,0) {mid-train};
    \node[modhi,   minimum width=15mm] (c) at (3.90,0) {SFT};
    \node[modhi,   minimum width=15mm] (d) at (5.85,0) {pref.\\opt.};
    \node[modhi,   minimum width=15mm] (e) at (7.80,0) {RLVR};
    \node[mod,     minimum width=15mm] (f) at (9.75,0) {distil.\\merge};
    \draw[flow] (a) -- (b);
    \draw[flow] (b) -- (c);
    \draw[flow] (c) -- (d);
    \draw[flow] (d) -- (e);
    \draw[flow] (e) -- (f);
    \node[lbl, below=5mm of a, align=center] {week 1};
    \node[lbl, below=5mm of f, align=center] {week 2};
  \end{tikzpicture}

  \vspace{0.6em}
  \raggedright
  \small
  \begin{tabular}{@{}L{2.2cm} L{8.9cm}@{}}
    \toprule
    \textbf{SFT} & supervised fine-tuning: teaches the format and \emph{a} good
      answer. Cannot teach ``better than'', and needs written responses\\ \addlinespace[0.3em]
    \textbf{pref.\ opt.} & teaches ``better than'' from pairwise human judgement.
      Bottlenecked by preference labels\\ \addlinespace[0.3em]
    \textbf{RLVR} & reinforcement learning with verifiable rewards: removes the
      human wherever the answer can be checked by a program, so it scales with
      compute rather than with the number of annotators\\ \addlinespace[0.3em]
    \textbf{distil.} & recovers what the RL stages damaged, and collapses many
      specialists into one shippable checkpoint\\
    \bottomrule
  \end{tabular}
  \note{This table is the spine of the next hour. Each stage exists because the
  previous one has a named limitation, and if they remember only the right-hand
  column they can reconstruct the lecture.

  Note where the boundary with week 1 sits, and say honestly that it is blurry.
  Mid-training uses the same objective and the same optimiser as pretraining,
  and teams disagree about where it ends.}
\end{frame}

% Optional reference: a fast-moving chronology and model-specific recipes.
\iffalse
\begin{frame}[shrink=10]{What the field was excited about, year by year}
  \centering
  \begin{tabular}{@{}c L{5.0cm} L{4.8cm}@{}}
    \toprule
    \textbf{Year} & \textbf{Dominant focus} & \textbf{Why}\\
    \midrule
    2022 & RLHF + PPO & produced the original ChatGPT\\
    2023 & LoRA SFT & small custom models on one GPU\\
    2024 & mid-training & synthetic data, data mixes, long context\\
    2025 & RLVR + GRPO & DeepSeek-R1; reasoning models\\
    2026 & RLVR beyond maths and code, plus inference-time scaling & verifiable
      environments for agents and computer use\\
    2027 & continual learning \cit{predicted} & the forgetting problem, taken
      seriously\\
    \bottomrule
  \end{tabular}

  \vspace{0.8em}
  \raggedright
  \small
  The list is \key{cumulative}. Nothing on it got replaced. The older methods
  just stopped being the part people wrote about, and every stage in the
  previous slide is still in every flagship recipe.
  \source{Raschka, \textit{The State of LLMs 2025} and
  \textit{Build Your Own Reasoning Model}, October 2026.}
  \note{The cumulative point is the one students get wrong. They read that DPO
  ``replaced'' RLHF and that GRPO ``replaced'' DPO. Neither is true: a 2026
  flagship model goes through SFT, some preference-style training and RLVR, in
  that order, and then a distillation pass.

  The 2027 row is a prediction and should be labelled as one out loud.}
\end{frame}

\begin{frame}[shrink=11]{Two real recipes, 2026}
  \begin{columns}[T]
    \begin{column}{0.48\textwidth}
      \begin{block}{Qwen3 --- four stages}
        \small
        \begin{enumerate}
          \item long chain-of-thought SFT
          \item reasoning RL
          \item \key{thinking-mode fusion} --- SFT on a mixture of thinking
            and non-thinking examples
          \item general RL for instruction and format following, in both modes
        \end{enumerate}
      \end{block}
    \end{column}
    \begin{column}{0.48\textwidth}
      \begin{block}{GLM-5 --- five stages}
        \small
        \begin{enumerate}
          \item multi-task SFT (interleaved, preserved and turn-level thinking)
          \item reasoning RL
          \item agentic RL
          \item general RL
          \item \key{on-policy self-distillation} from the earlier checkpoints
        \end{enumerate}
      \end{block}
    \end{column}
  \end{columns}
  \vspace{0.5em}
  \raggedright
  Also worth knowing: \textbf{DeepSeek V4} runs a separate RL configuration per
  effort mode, each with its own context window and its own length penalty, and
  then distils more than ten domain specialists plus three effort specialists
  into one checkpoint. \textbf{Kimi K3} does the same with nine specialists.
  \source{Raschka, \textit{Controlling Reasoning Effort in LLMs}, July 2026,
  summarising the Qwen3, GLM-5, DeepSeek V4 and Kimi K3 technical reports.}
  \note{Pick out the shared shape: supervised stage first to install the
  format, then reinforcement learning in increasingly specific domains, then a
  cleanup pass.

  Say the quiet part too. Nobody publishes the ablations that would tell you
  which of the five stages mattered, so what we are reading here is a set of
  recipes rather than a set of experiments.}
\end{frame}
\fi

\begin{frame}{The pipeline forgets its own earlier stages}
  \begin{columns}[T]
    \begin{column}{0.56\textwidth}
      GLM-5's final distillation step exists, in the authors' own words,
      \emph{to recover capabilities that may have weakened during the
      sequential RL stages}.
      \vspace{0.6em}

      That is \key{catastrophic forgetting inside the post-training pipeline
      itself}. Four RL stages aimed at reasoning, agents and general behaviour
      each degraded something the earlier stages had installed, and the fix was
      to distil from the checkpoints that still had it.
      \vspace{0.6em}

      You will run into the same thing on Thursday with a single GPU.
    \end{column}
    \begin{column}{0.44\textwidth}
      \begin{takeaway}
        \small
        Every post-training stage is a trade. Asking whether it improved the
        target metric is not enough; you also have to ask what it cost
        somewhere else.
        \vspace{0.4em}

        Lab 2 measures exactly this, with a fifty-question general-knowledge
        check before and after your fine-tune.
      \end{takeaway}
    \end{column}
  \end{columns}
  \source{Raschka, \textit{Controlling Reasoning Effort in LLMs}, July 2026,
  quoting the GLM-5 technical report.}
  \note{This is the slide that justifies the third of the three things to
  remember, so deliver it deliberately.

  If someone asks why they cannot just mix the old data back in: that is
  exactly what rehearsal and self-distillation are, and it is what GLM-5 did.
  The price is a whole extra training stage.

  Clock check: should be around 16:30 here. Section 3 is fifteen minutes and
  it is the one they need on Thursday.}
\end{frame}

% =====================================================================
\section{Supervised fine-tuning}
% =====================================================================

\begin{frame}{SFT is pretraining with better manners}
  \[
    \mathcal{L}_{\mathrm{SFT}} \;=\;
    -\sum_{t=1}^{T} m_t \,\log p_\theta\!\left(x_t \mid x_{<t}\right)
    \qquad\text{with } m_t \in \{0,1\}
  \]
  \emph{In words:} the week 1 next-token loss, with an on/off switch $m_t$ in
  front of every token. Switch it off and that token is read but not learned.
  \vspace{0.2em}

  Compared with week 1, exactly three things change:
  \vspace{0.4em}
  \begin{enumerate}
    \item the \textbf{data} is conversations rendered through the model's own
      chat template, not web documents
    \item the \textbf{mask} $m_t$ zeroes the loss on everything that is not
      assistant content
    \item the \textbf{budget} is tiny: thousands of examples and one GPU,
      against trillions of tokens and a cluster
  \end{enumerate}
  \vspace{0.5em}
  Nothing happens to the optimiser, the objective or the architecture. That is
  what makes SFT cheap, and it is also the reason it cannot do very much.
  \note{Start by writing the loss and pointing at the only new symbol. Students
  who think fine-tuning is a different algorithm relax when they see it is the
  same cross-entropy with an indicator in front.

  The third point is the honest framing of the whole lecture. We are nudging a
  distribution that was shaped by ten trillion tokens, and we are doing it with
  a few million. Expecting that to add knowledge would be strange.}
\end{frame}

\begin{frame}{The loss mask, drawn}
  \centering
  \begin{tikzpicture}[x=0.84cm]
    % group labels
    \node[lbl] at (0.5,0.62)  {template};
    \node[lbl] at (3.5,0.62)  {user content};
    \node[lbl] at (7.0,0.62)  {template};
    \node[lbl] at (10.5,0.62) {assistant content};
    % token cells
    \node[tokmask, minimum width=0.80cm] at (0,0)  {\texttt{im\_s}};
    \node[tokmask, minimum width=0.80cm] at (1,0)  {\texttt{user}};
    \node[tokoff,  minimum width=0.80cm] at (2,0)  {\texttt{Tell}};
    \node[tokoff,  minimum width=0.80cm] at (3,0)  {\texttt{me}};
    \node[tokoff,  minimum width=0.80cm] at (4,0)  {\texttt{a}};
    \node[tokoff,  minimum width=0.80cm] at (5,0)  {\texttt{joke}};
    \node[tokmask, minimum width=0.80cm] at (6,0)  {\texttt{im\_e}};
    \node[tokmask, minimum width=0.80cm] at (7,0)  {\texttt{im\_s}};
    \node[tokmask, minimum width=0.80cm] at (8,0)  {\texttt{asst}};
    \node[tokon,   minimum width=0.80cm] at (9,0)  {\texttt{Sure}};
    \node[tokon,   minimum width=0.80cm] at (10,0) {\texttt{here}};
    \node[tokon,   minimum width=0.80cm] at (11,0) {\texttt{goes}};
    \node[tokon,   minimum width=0.80cm] at (12,0) {\texttt{im\_e}};
    % mask row
    \node[lbl] at (-1.25,-0.62) {$m_t$};
    \node[lbl] at (0,-0.62)  {0};
    \node[lbl] at (1,-0.62)  {0};
    \node[lbl] at (2,-0.62)  {0};
    \node[lbl] at (3,-0.62)  {0};
    \node[lbl] at (4,-0.62)  {0};
    \node[lbl] at (5,-0.62)  {0};
    \node[lbl] at (6,-0.62)  {0};
    \node[lbl] at (7,-0.62)  {0};
    \node[lbl] at (8,-0.62)  {0};
    \node[lbl] at (9,-0.62)  {1};
    \node[lbl] at (10,-0.62) {1};
    \node[lbl] at (11,-0.62) {1};
    \node[lbl] at (12,-0.62) {1};
  \end{tikzpicture}

  \vspace{0.5em}
  \raggedright
  \small
  Abbreviations: \texttt{im\_s} is
  \texttt{\textless\textbar im\_start\textbar\textgreater},
  \texttt{im\_e} is
  \texttt{\textless\textbar im\_end\textbar\textgreater},
  \texttt{asst} is \texttt{assistant}.
  \vspace{0.3em}

  All of it is still \emph{in the context}, and the model attends to all of it.
  Only the blue cells contribute a gradient. Notice that the final end-of-turn
  token sits \key{inside} the mask, because stopping is a behaviour you have to
  train.
  \note{Draw the four regions with your hand as you talk. The distinction
  people miss is between \emph{being in the context} and \emph{being in the
  loss}, so say that sentence explicitly.

  The last line is the one that saves them debugging time. If a fine-tuned
  model rambles past the end of its answer, the first thing to check is whether
  the end-of-turn token was masked out of the loss.}
\end{frame}

\begin{frame}{What goes wrong if you do not mask}
  \begin{columns}[T]
    \begin{column}{0.5\textwidth}
      \textbf{Train on everything and the model learns to be the user too.}
      \begin{itemize}
        \item it will cheerfully generate the next \emph{user} turn after
          finishing its own
        \item the gradient is spent reproducing text you already have
        \item with long prompts and short answers (classification, NER,
          extractive QA) most of your compute goes on memorising inputs
        \item the loss curve looks beautiful, because predicting your own
          prompts is easy
      \end{itemize}
    \end{column}
    \begin{column}{0.5\textwidth}
      \begin{block}{In practice, on Thursday}
        \small
        TRL's \texttt{SFTTrainer} will do this for you, but only if you ask:
        \texttt{assistant\_only\_loss=True}.
        \vspace{0.4em}

        That flag requires the chat template to mark the assistant region.
        If the template does not, you get the whole sequence in the loss and
        \textbf{no warning}.
        \vspace{0.4em}

        Check it by printing the labels of one batch and confirming that the
        ignored positions are exactly the prompt.
      \end{block}
    \end{column}
  \end{columns}
  \vspace{0.5em}
  \raggedright
  \small
  Honest caveat: training on prompts is not catastrophic, and a few recipes do
  it deliberately as a regulariser. You just do not want it when the prompt is
  long and the answer is three tokens.
  \note{The ``print one batch and look at the labels'' instruction is the most
  practically valuable sentence in the section. Tell them to do it once, for
  real, on Thursday, before they start a long run.

  The caveat at the bottom protects you from the student who has read a paper
  where prompt loss was kept. The answer is that it depends on the ratio of
  prompt length to answer length, and for their tasks that ratio is bad.}
\end{frame}

\begin{frame}[shrink=4]{Data: a thousand good examples go a long way}
  \begin{columns}[T]
    \begin{column}{0.54\textwidth}
      The argument, in one line: SFT \key{selects} a behaviour the model
      already has, and to do that you need clean, consistent demonstrations
      rather than a lot of them.
      \vspace{0.5em}
      \begin{itemize}
        \item LIMA fine-tuned a 65B model on \textbf{1{,}000} curated
          examples and remained competitive with far larger instruction sets
          \cit{Zhou et al.\ 2023}
        \item one inconsistent example teaches the model that the format is
          optional
        \item a scraped set teaches you the average of its annotators,
          including their worst day
      \end{itemize}
    \end{column}
    \begin{column}{0.46\textwidth}
      \begin{block}{What ``quality'' concretely means}
        \small
        \textbf{consistent format} --- the same answer shape every time

        \textbf{correct} --- you checked, with a script where possible

        \textbf{on-distribution} --- the prompts look like the prompts you will
        receive in practice

        \textbf{diverse} --- across the input space, not across phrasings of
        the same input

        \textbf{deduplicated} --- against each other and against your test set
      \end{block}
    \end{column}
  \end{columns}
  \vspace{0.4em}
  \raggedright
  \small
  For Lab 2, \textbf{1--5k examples} is the right order of magnitude. You will
  get more out of an afternoon cleaning them than out of a week collecting
  more.
  \note{Flag the sourcing honestly: the LIMA result is from the standard
  literature, not from the Raschka corpus the rest of this lecture leans on.
  Modelling that distinction is part of the course.

  The decontamination point in the last bullet is the one that bites them in
  Lab 4. If their SFT set overlaps their held-out set, every number they report
  for the rest of term is wrong.}
\end{frame}

% Optional catalogue; the core lecture focuses on data quality and masking.
\iffalse
\begin{frame}[shrink=19]{Where SFT data comes from}
  \centering
  \begin{tabular}{@{}L{2.6cm} L{4.2cm} L{4.2cm}@{}}
    \toprule
    \textbf{Source} & \textbf{Good for} & \textbf{Costs and risks}\\
    \midrule
    human-written & tone, safety, genuinely novel formats & slow, expensive,
      inconsistent between annotators\\[0.3em]
    converted from existing labelled data & cheap, exactly on-task; what you
      will do on Thursday & teaches terse labels, not explanations; can look
      unnatural\\[0.3em]
    generated by a larger model (distillation) & scales to any volume, gives
      you reasoning traces you could not write & inherits the teacher's errors
      and style; check the teacher's licence\\[0.3em]
    generated then filtered by a verifier & the best of both where the task is
      checkable & needs a verifier, which is section 5\\
    \bottomrule
  \end{tabular}

  \vspace{0.7em}
  \raggedright
  \small
  Raschka's worked example of row three: generate answers to the MATH dataset
  with \textbf{DeepSeek-R1} as the teacher, train a small model on them, and
  evaluate on MATH-500. He presents it explicitly as \emph{an alternative to
  developing reasoning capability through reinforcement learning}.
  \source{Raschka, \textit{Build Your Own Reasoning Model}, October 2026.}
  \note{Row two is what most of the room will do: take their task's training
  split and render it as conversations. Warn them that a dataset of one-word
  gold labels produces a model that only ever emits one word, which is fine
  for accuracy and useless for anything else.

  Row three is the optional distillation exercise in Lab 2, and it is also a
  licence question. Some model licences restrict training other models on their
  outputs. Tell them to read it, because the report asks for licences.}
\end{frame}

% Distillation taxonomy is also kept inside this optional block.
\begin{frame}{Distillation in 2026: many teachers feeding one student}
  \begin{columns}[T]
    \begin{column}{0.46\textwidth}
      \textbf{The old picture (2023)}\\
      one big teacher $\rightarrow$ one small student, and the student ends up
      as a cheaper copy.
      \vspace{0.8em}

      \textbf{The 2026 picture}\\
      many \emph{specialist} teachers $\rightarrow$ one \emph{unified} student,
      which ends up being something none of the teachers were.
    \end{column}
    \begin{column}{0.54\textwidth}
      \begin{block}{Concretely}
        \small
        \textbf{DeepSeek V4} --- more than ten domain specialists plus three
        effort specialists (Non-think, Think High, Think Max), distilled
        on-policy into one checkpoint

        \textbf{Kimi K3} --- nine specialists, three effort levels across three
        domains, multi-teacher on-policy distillation

        \textbf{GLM-5} --- teachers are its \emph{own} earlier checkpoints

        \textbf{Nemotron 3 Ultra} --- learns its medium-effort mode from
        gpt-oss-120B traces generated in medium-effort mode
      \end{block}
    \end{column}
  \end{columns}
  \vspace{0.4em}
  \raggedright
  \small
  This is how multi-mode models get built. Rather than training one model to do
  five things, you train five models and merge them by imitation.
  \source{Raschka, \textit{Controlling Reasoning Effort in LLMs}, July 2026.
  Model facts as of October 2026.}
  \note{The Nemotron row is a nice detail: a closed-ish capability, effort
  control, was transferred from one open-weight model to another just by
  sampling it in the right mode. That one example covers most of the economics
  of distillation.

  Ask the room what the risk is. The answer you want is homogenisation: if
  everyone distils from the same three teachers, the apparent diversity of the
  open-weight field is thinner than it looks.}
\end{frame}

\begin{frame}[shrink=3]{The distillation $2\times2$}
  \centering
  \begin{tabular}{@{}L{2.2cm} L{4.3cm} L{4.3cm}@{}}
    \toprule
    & \textbf{off-policy} \newline (teacher's own outputs) &
      \textbf{on-policy} \newline (student's outputs, teacher-scored)\\
    \midrule
    \textbf{hard} \newline (sampled tokens) &
      the classic recipe: collect teacher answers, run SFT on them &
      student generates, teacher or verifier keeps the good ones; rejection
      sampling\\[0.4em]
    \textbf{soft} \newline (full distributions) &
      match the teacher's next-token distribution over its own text &
      match it \emph{on the student's own trajectories}; the 2026 flagships'
      ``on-policy distillation''\\
    \bottomrule
  \end{tabular}

  \vspace{0.6em}
  \raggedright
  \small
  On-policy works better because the student gets corrected in the states it
  actually reaches. It is rarer because you have to run the teacher throughout
  training. Soft distillation also needs \textbf{the same tokenizer}, which
  rules out most cross-family pairs.
  \source{Taxonomy from Raschka, \textit{Build Your Own Reasoning Model},
  October 2026.}
  \note{Spell out the on-policy argument a little more than the slide does. An
  off-policy student has to generalise from the teacher's distribution to its
  own, and an on-policy student never has to, because it is being scored on
  text it generated itself. The price is that the bottom-right cell is the
  expensive quadrant: the teacher has to be in memory and running for the whole
  of training.

  The tokenizer constraint is the practical one and it is rarely stated. If the
  teacher and student do not share a vocabulary you cannot match distributions,
  so you are restricted to hard distillation.

  For Lab 2, everything they do is the top-left cell. Say that, so the table
  does not feel like homework.}
\end{frame}
\fi

\begin{frame}[shrink=4]{Catastrophic forgetting}
  \begin{columns}[T]
    \begin{column}{0.54\textwidth}
      \textbf{What it is.} Learning new behaviour moves weights, and the
      behaviour those weights previously encoded degrades. In continued
      pretraining this is well documented; in SFT it is the same mechanism
      with a smaller step size.
      \vspace{0.5em}

      \textbf{Why it is unavoidable.} There is no free capacity lying around.
      A gradient computed on a narrow distribution overwrites a representation
      that the rest of the model was sharing.
      \vspace{0.5em}

      \textbf{What makes it worse:} high learning rate, many epochs, narrow
      data, full fine-tuning rather than LoRA, and a small model.
    \end{column}
    \begin{column}{0.46\textwidth}
      \begin{block}{What it looks like}
        \small
        Your task metric goes up by 12 points.

        The model now answers every question in your task's format, including
        ``what is the capital of Finland''.

        It has stopped using the second language it used to handle.

        It no longer refuses anything.

        \textit{None of this appears in your task metric.}
      \end{block}
    \end{column}
  \end{columns}
  \vspace{0.3em}
  \raggedright
  \small
  Mitigations, in increasing cost: fewer epochs \textbullet\ lower learning
  rate \textbullet\ LoRA with a small rank \textbullet\ mixing general data
  back in (rehearsal) \textbullet\ distilling from your own pre-SFT checkpoint.
  \note{The right-hand box is deliberately concrete. Students expect forgetting
  to look like amnesia. In practice it usually looks like the model applying
  its new format to everything.

  Note that the last two mitigations are exactly what GLM-5 did at industrial
  scale, with the same techniques they will use on one GPU.}
\end{frame}

\begin{frame}{Measure what broke, not only what improved}
  \begin{columns}[T]
    \begin{column}{0.55\textwidth}
      Use a multiple-choice benchmark such as \textbf{MMLU} as a
      \key{regression test}, not as a capability score.
      \vspace{0.5em}
      \begin{itemize}
        \item score the base model, score the fine-tuned model, report the
          \emph{difference}
        \item the absolute number tells you very little. The change tells you
          a lot
        \item a high score is weak evidence that the model is good, but a drop
          is strong evidence that you broke something
      \end{itemize}
      \vspace{0.5em}
      Random guessing on four options is \textbf{25\%}, so anything near that
      is a scoring bug, not a result.
    \end{column}
    \begin{column}{0.45\textwidth}
      \begin{takeaway}
        \small
        In Lab 2 this is the fourth column of your results table: fifty
        general-knowledge questions, scored \emph{before} and \emph{after}, with
        identical decoding settings.
        \vspace{0.4em}

        A report with three numbers and no forgetting check is incomplete.
      \end{takeaway}
    \end{column}
  \end{columns}
  \source{Raschka, \textit{Understanding the 4 Main Approaches to LLM
  Evaluation}: multiple-choice benchmarks are most useful for
  ``checking if and how much knowledge the model has forgotten compared to
  the base model''.}
  \note{This is the third of the three things to remember, introduced here and
  repeated at the end. Say the sentence in the form you want it quoted back:
  always measure what fine-tuning broke, not only what it improved.

  The identical-decoding-settings instruction matters more than it sounds. If
  the two runs used different temperatures, you have measured the temperature
  and not the fine-tune.}
\end{frame}

% Optional advanced example of loss-mask design.
\iffalse
\begin{frame}{A more deliberate use of the loss mask}
  \textbf{Nemotron 3 Ultra: training a model to be interrupted.}
  \vspace{0.5em}
  \begin{columns}[T]
    \begin{column}{0.56\textwidth}
      The recipe:
      \begin{enumerate}
        \item take ordinary long reasoning traces
        \item \textbf{truncate each at a random token budget}
        \item inject a closing \texttt{\textless/think\textgreater} token
        \item keep the original final answer
        \item \key{mask the injected token out of the SFT loss}
      \end{enumerate}
    \end{column}
    \begin{column}{0.44\textwidth}
      \begin{block}{Why the mask}
        \small
        Leave the injected token in the loss and the model learns to
        \emph{emit} it, which means cutting its own reasoning short at random.
        \vspace{0.4em}

        Mask it and the model only learns to \emph{react} when something
        external inserts it: stop thinking, produce the best answer you have
        so far.
      \end{block}
    \end{column}
  \end{columns}
  \vspace{0.5em}
  \raggedright
  The result is a model that degrades gracefully under an externally imposed
  token budget, which is what makes a hard inference-time budget usable.
  \source{Raschka, \textit{Controlling Reasoning Effort in LLMs}, July 2026.}
  \note{This is the most precise use of a loss mask in the corpus and it is
  worth ninety seconds, because it shows masking as a design tool rather than a
  default setting. Masking is how you get ``learn to respond when X happens''
  instead of ``learn to do X''.

  Related and worth mentioning: Qwen3 can also be stopped mid-thought, and its
  technical report says that behaviour was never explicitly trained. It emerged
  from thinking-mode fusion.

  Clock check: should be around 16:45 here. Section 4 starts now.}
\end{frame}
\fi

% =====================================================================
\section{Preference optimisation}
% =====================================================================

\begin{frame}[shrink=4]{Why SFT is not enough}
  \begin{columns}[T]
    \begin{column}{0.54\textwidth}
      SFT maximises the likelihood of \key{one} demonstrated answer. It has
      no way to express that one answer is \key{better than} another.
      \vspace{0.6em}

      That matters whenever the target is a \emph{judgement} rather than a
      fact:
      \begin{itemize}
        \item helpfulness, tone, length, hedging
        \item refusing the right things and only those
        \item ``answer the question that was asked''
        \item formatting that nobody can write down as a rule
      \end{itemize}
      \vspace{0.4em}
      Writing the gold answer for any of these is hard, while \emph{choosing}
      between two candidates is easy, and every method in this section is
      built on that gap.
    \end{column}
    \begin{column}{0.46\textwidth}
      \begin{block}{The second problem SFT has}
        \small
        SFT never sees a bad answer. It is trained only on positives, so it
        has no gradient that pushes probability mass \emph{away} from the
        failure modes it will actually produce.
        \vspace{0.4em}

        Every method in this section is, mechanically, a way to add that
        negative signal.
      \end{block}
    \end{column}
  \end{columns}
  \note{Judging is easier than generating. That is why this whole family of
  methods exists, and it comes back in week 4 as the justification for
  LLM-as-a-judge.

  The right-hand box is the sharper technical statement to remember: SFT has
  only positive examples.}
\end{frame}

\begin{frame}{RLHF in one diagram}
  \centering
  \begin{tikzpicture}
    \node[mod, minimum width=20mm] (p)  at (0,0)   {prompts};
    \node[mod, minimum width=22mm] (s)  at (2.6,0) {sample 2--$k$\\responses};
    \node[mod, minimum width=22mm] (h)  at (5.4,0) {humans rank\\them};
    \node[modhi, minimum width=22mm] (r) at (8.2,0) {reward model\\$r_\phi(x,y)$};
    \node[modhi, minimum width=26mm] (ppo) at (5.4,-1.9) {PPO: maximise $r_\phi$\\under a KL leash};
    \node[modmute, minimum width=22mm] (ref) at (1.5,-1.9) {frozen reference\\$\pi_{\mathrm{ref}}$ (the SFT model)};
    \draw[flow] (p) -- (s);
    \draw[flow] (s) -- (h);
    \draw[flow] (h) -- (r);
    \draw[flow] (r.south) |- (ppo.east);
    \draw[flow] (ref) -- (ppo);
  \end{tikzpicture}

  \vspace{0.5em}
  \raggedright
  \small
  The RL vocabulary, once, in plain words. The \textbf{policy} is the model
  being trained. The \textbf{reference} is a frozen copy of it from before this
  stage. The \textbf{KL leash} is a penalty for drifting away from that copy.
  The \textbf{value model} (or critic) is a second network whose only job is to
  guess how good an answer will turn out.
  \vspace{0.2em}

  That is four large models in memory at once, which is why the field went
  looking for something simpler.
  \source{Ouyang et al.\ 2022, \textit{Training language models to follow
  instructions with human feedback} (InstructGPT).}
  \note{Walk the diagram left to right once, then point at the memory
  accounting underneath. The intuition to leave them with: the algorithm is not
  exotic, and what makes classic RLHF expensive is holding three or four copies
  of a large model while generating from one of them inside the training loop.

  The KL leash deserves a sentence: without it the policy walks off and
  produces degenerate text that the reward model happens to like.}
\end{frame}

\begin{frame}{The reward model is logistic regression on pairs}
  Collect triples: a prompt $x$, the answer a human \emph{preferred} ($y_w$)
  and the one they \emph{rejected} ($y_l$). Train a scorer $r_\phi$ so that
  \[
    P(y_w \text{ wins}) \;=\;
    \sigma\!\left(r_\phi(x,y_w) - r_\phi(x,y_l)\right)
  \]
  \emph{In words:} the bigger the score gap between two answers, the more
  confidently we predict which one a human picked. You fit it like any logistic
  regression, by maximising the probability it assigns to the choices people
  made. The recipe has a name, \key{Bradley--Terry}, and it returns in week 4
  for arena leaderboards.
  \vspace{0.4em}

  \begin{columns}[T]
    \begin{column}{0.5\textwidth}
      \small
      $r_\phi$ is usually the SFT model with its next-token head swapped for a
      head that emits a single number. Only \emph{differences} of scores are
      pinned down, because the scale has no absolute zero.
    \end{column}
    \begin{column}{0.5\textwidth}
      \small
      The RL stage then looks for the model that earns the highest average
      score from $r_\phi$, minus a penalty $\beta$ times how far it has drifted
      from the frozen starting model.
    \end{column}
  \end{columns}
  \note{Keep this at the level of ``it is logistic regression on pairs''. The
  detail that matters later is that only differences are identified, since that
  is the degree of freedom DPO exploits when it absorbs the partition function.

  The two expressions taken off the slide, for the board or for a question.
  The reward model is fitted by minimising
  $\mathcal{L}_{\mathrm{RM}} = -\mathbb{E}_{(x,y_w,y_l)\sim\mathcal{D}}
  [\log \sigma(r_\phi(x,y_w) - r_\phi(x,y_l))]$, which is exactly the
  cross-entropy of a two-class logistic regression whose single feature is the
  score difference. The RL objective is then
  $\max_{\pi_\theta} \mathbb{E}_{y\sim\pi_\theta}[r_\phi(x,y)]
  - \beta\,\mathbb{D}_{\mathrm{KL}}[\pi_\theta \| \pi_{\mathrm{ref}}]$, where
  the KL term measures, per token, how much the trained model's probability
  distribution has moved away from the frozen one.

  If anyone asks why not just regress on absolute scores: human absolute
  ratings are famously unstable, pairwise comparisons are not.}
\end{frame}

\begin{frame}{PPO, briefly, because GRPO is a simplification of it}
  \[
    \mathcal{J}_{\mathrm{PPO}} = \mathbb{E}\!\left[
      \min\!\left(\rho \hat{A},\;
      \mathrm{clip}(\rho)\,\hat{A}\right)\right],
    \qquad
    \rho = \frac{\pi_{\mathrm{new}}(\text{this token})}
                {\pi_{\mathrm{old}}(\text{this token})}
  \]
  \emph{In words:} make the tokens of better-than-average answers more likely
  and the rest less likely, with a cap on how far one batch may move the model.
  \vspace{0.3em}

  \begin{columns}[T]
    \begin{column}{0.52\textwidth}
      \small
      \textbf{Advantage} $\hat{A}$ --- how much better than average this answer
      turned out to be. ``Average'' has to come from somewhere.
      \vspace{0.4em}

      \textbf{Clipping} --- the cap. The ratio is held within about 20\% of 1.
      Without it, one bad batch wrecks the model.
      \vspace{0.4em}

      \textbf{The critic} --- a second network, as large as the one being
      trained, whose job is to predict that average.
    \end{column}
    \begin{column}{0.48\textwidth}
      \begin{block}{The part that matters next}
        \small
        The critic exists only to supply that ``average'', and it is another
        model the size of the one being trained, with its own loss and its own
        failure modes.
        \vspace{0.4em}

        \key{GRPO's contribution is a cheaper way to get the same average.}
      \end{block}
    \end{column}
  \end{columns}
  \source{Schulman et al.\ 2017, \textit{Proximal Policy Optimization
  Algorithms}.}
  \note{Do not try to teach PPO properly. Teach the one thing section 5 needs,
  which is the role of the critic, and treat the rest as scaffolding.

  Full notation, if someone wants it. The expectation runs over time steps $t$;
  the ratio is
  $\rho_t = \pi_\theta(a_t\mid s_t)/\pi_{\theta_{\mathrm{old}}}(a_t\mid s_t)$,
  where $a_t$ is the token emitted at step $t$ and $s_t$ is everything before
  it; and the clip is the two-sided
  $\mathrm{clip}(\rho_t, 1-\epsilon, 1+\epsilon)$. Taking the minimum of the
  clipped and unclipped terms makes the objective pessimistic: the cap only
  ever bites in the direction that would have helped.

  The clipped objective reappears almost verbatim in GRPO, so flag it now:
  when they see the GRPO loss in fifteen minutes, the only genuinely new part
  will be how the advantage is computed.}
\end{frame}

\begin{frame}{Reward source and optimisation are separate choices}
  \centering
  \small
  \begin{tabular}{@{}L{2.6cm} L{4.1cm} C{1.5cm} C{1.5cm}@{}}
    \toprule
    & \textbf{Reward source} & \textbf{Reward model?} & \textbf{Critic?}\\
    \midrule
    \textbf{RLHF + PPO} & learned human-preference score & yes & yes\\
    \textbf{GRPO} & learned score & yes & \bad{no}\\
    \textbf{RLVR + GRPO} & programmatic verifier & \bad{no} & \bad{no}\\
    \bottomrule
  \end{tabular}

  \vspace{0.3em}
  \raggedright\small
  \begin{itemize}
    \item \textbf{GRPO removes the critic}: it compares several responses to
      the same prompt instead of predicting a baseline with another network.
    \item \textbf{RLVR changes the reward source}: a program checks the answer.
      GRPO can also be used with a learned reward.
    \item \textbf{DPO takes a different route}: it trains on stored preference
      pairs, without generating responses during the update.
  \end{itemize}
  \note{Put this table up and leave it up for a beat. It is one of the most
  reusable objects in the lecture.

  The ``two separate deletions'' framing is what stops students thinking GRPO
  and RLVR are the same thing. GRPO is an optimiser, whereas RLVR is a claim
  about where the reward comes from. You can run GRPO with a learned reward
  model, and people do.}
\end{frame}

\begin{frame}{DPO learns directly from preferred and rejected answers}
  A DPO example contains a prompt, a \textbf{preferred answer} and a
  \textbf{rejected answer}. Training increases the preferred answer's
  likelihood relative to the rejected one.
  \vspace{0.5em}

  \begin{columns}[T]
    \begin{column}{0.5\textwidth}
      \textbf{Why it is useful}
      \begin{itemize}
        \item people often find comparison easier than writing an ideal answer;
        \item training resembles ordinary supervised learning;
        \item there is no separate reward model or rollout loop.
      \end{itemize}
    \end{column}
    \begin{column}{0.5\textwidth}
      \textbf{What to watch}
      \begin{itemize}
        \item the pair only tells us which answer is better;
        \item the data may reflect annotator or judge bias;
        \item a frozen reference model keeps the update close to its starting
          behaviour.
      \end{itemize}
    \end{column}
  \end{columns}

  \begin{takeaway}
    Use DPO when you have reliable preference pairs and the target quality is
    difficult to express as an automatic score.
  \end{takeaway}
  \source{Rafailov et al.\ 2023, \textit{Direct Preference Optimization: Your
  Language Model is Secretly a Reward Model}.}
  \note{Keep the projected explanation conceptual. The mathematical derivation
  and loss are retained in the source as optional reading. Connect DPO to the
  earlier reward-model slide: both learn from comparisons, but DPO updates the
  language model directly.}
\end{frame}

% The derivation and variant catalogue are reading material; the projected
% lecture keeps the idea and its main limitation.
\iffalse
\begin{frame}[shrink=12]{The DPO loss}
  \[
    \mathcal{L}_{\mathrm{DPO}} = -\,\mathbb{E}
    \left[\log \sigma\!\left(
      \beta \log\frac{\pi_\theta(y_w\mid x)}{\pi_{\mathrm{ref}}(y_w\mid x)}
      \;-\;
      \beta \log\frac{\pi_\theta(y_l\mid x)}{\pi_{\mathrm{ref}}(y_l\mid x)}
    \right)\right]
  \]
  \emph{In words:} for each preference pair, see how much more likely the
  trained model makes the \emph{preferred} answer than the frozen copy does,
  do the same for the \emph{rejected} answer, and push those two apart.
  \vspace{0.2em}

  \centering
  \small
  \begin{tabular}{@{}L{2.6cm} L{8.5cm}@{}}
    \toprule
    $\pi_\theta$, $\pi_{\mathrm{ref}}$ & the model being trained, and a frozen
      copy of the SFT model it started from. The frozen copy is what
      ``unchanged'' means. It costs one extra forward pass per batch, with no
      gradients and no optimiser state\\ \addlinespace[0.25em]
    $\beta$ & how hard the leash pulls. Large (0.5): stay close, learn little.
      Small (0.01): drift freely, lose fluency. Typical: \textbf{0.1}\\ \addlinespace[0.25em]
    $\sigma$ & the logistic function: it turns that gap into a probability.
      Big gap, small loss\\
    \bottomrule
  \end{tabular}

  \vspace{0.6em}
  \raggedright
  No reward model appears anywhere in this expression, and no text is generated
  during training. DPO costs about as much as SFT plus one frozen forward pass.
  \note{The row on the reference model is the one that costs people memory on
  Thursday. If they try Bonus A, remind them that DPO needs the reference in
  memory as well as the policy, and that with LoRA you can get it for free by
  disabling the adapter rather than loading a second copy.

  Term by term, if you want to walk it on the board. The expectation runs over
  triples $(x, y_w, y_l)$ drawn from the preference dataset $\mathcal{D}$;
  $y_w$ is the ``winning'' (chosen) response and $y_l$ the ``losing''
  (rejected) one. Each log-ratio
  $\beta\log \pi_\theta(y\mid x)/\pi_{\mathrm{ref}}(y\mid x)$ is exactly the
  implicit reward from the previous slide, with the partition term already
  cancelled. Their difference is the \emph{margin}; $\log\sigma$ of the margin
  is the log-likelihood of a Bradley--Terry model that says the chosen response
  wins. So DPO is literally the reward-model loss, with the reward replaced by
  the model's own log-ratio.

  Beta is the only hyperparameter they will actually tune. Give the three
  numbers on the slide and tell them to start at 0.1.}
\end{frame}

\begin{frame}{DPO's known failure mode}
  \begin{gotcha}[Both likelihoods can fall]
    The loss only constrains the \textbf{gap} between the two log-ratios.
    Nothing in it requires the chosen answer to actually become more likely.
    \vspace{0.3em}

    In practice the log-probability of the \textbf{chosen} response and of the
    \textbf{rejected} response often both decrease during training, with the
    rejected one falling faster. Your margin improves while the model gets
    less and less likely to produce either answer.
  \end{gotcha}
  \vspace{0.4em}
  \begin{columns}[T]
    \begin{column}{0.5\textwidth}
      \small
      \textbf{Why it happens.} Probability mass has to go somewhere, and
      ``somewhere'' is unconstrained: off-distribution continuations neither
      annotator considered.
    \end{column}
    \begin{column}{0.5\textwidth}
      \small
      \textbf{What to do.} Log both log-probabilities, not just the loss and
      the implicit-reward accuracy. Add an SFT term on $y_w$ if they fall
      together. Reduce the learning rate before blaming the data.
    \end{column}
  \end{columns}
  \vspace{0.4em}
  \raggedright
  \small
  The deeper reason: DPO is \key{off-policy}. Some other model produced the
  pairs, so the training signal never covers the text your model is about to
  generate.
  \note{Gotcha two of three, and a useful piece of debugging advice. If they
  run Bonus A and the loss looks healthy while the outputs get worse, this is
  why.

  The off-policy point is the bridge to section 5. GRPO trains on the model's
  own rollouts, so it is on-policy by construction, and that difference turns
  out to matter.}
\end{frame}

\begin{frame}[shrink=13]{The relatives, one line each}
  \centering
  \begin{tabular}{@{}L{2.0cm} L{9.1cm}@{}}
    \toprule
    \textbf{IPO} & swaps the logistic loss for a squared-margin objective, so
      the optimum is finite rather than pushed to infinity. Fixes DPO's
      tendency to overfit deterministic preferences\\[0.35em]
    \textbf{KTO} & drops pairs entirely. Needs only \emph{unpaired} examples
      labelled desirable or undesirable, weighted by a prospect-theoretic
      utility \cit{Ethayarajh et al.\ 2024}\\[0.35em]
    \textbf{ORPO} & \textbf{no reference model at all}: an odds-ratio penalty
      added straight to the SFT loss, so alignment happens in a single stage
      \cit{Hong et al.\ 2024}\\[0.35em]
    \textbf{SimPO} & no reference model either. Uses the
      \emph{length-normalised} average log-probability as the implicit reward,
      plus a target margin, which removes DPO's length bias
      \cit{Meng et al.\ 2024}\\
    \bottomrule
  \end{tabular}

  \vspace{0.5em}
  \raggedright
  \small
  Notice the direction of travel. Each variant deletes a component: the
  pairing, the reference model, the length dependence. All four are in TRL's
  stable line, so trying one on Thursday is a configuration change.
  \note{Teach the pattern rather than four separate objectives. Every one of
  these is ``DPO minus something'', and knowing which component was deleted
  tells you what the method is for.

  KTO is the one with real production relevance, and it is also much cheaper
  to annotate for. Thumbs-up and thumbs-down logs are unpaired by nature, so
  insisting on pairs throws most of that data away.}
\end{frame}
\fi

\begin{frame}{Where preference machinery survives in 2026}
  \begin{columns}[T]
    \begin{column}{0.54\textwidth}
      It did not die. What happened is that it moved over to the
      \key{non-verifiable half} of the reward.
      \vspace{0.5em}

      \textbf{DeepSeek V3.2}, verbatim: rule-based outcome reward, length
      penalty and language-consistency reward \emph{for reasoning and agent
      tasks}; a \textbf{generative reward model where each prompt has its own
      rubrics} \emph{for general tasks}.
      \vspace{0.5em}

      So the 2026 pipeline is a \textbf{hybrid}: verifiers where the answer can
      be checked, LLM-as-a-judge with rubrics everywhere else.
      \textbf{Kimi K3} likewise keeps learned reward models and pairwise
      comparisons for its non-verifiable and agentic tasks.
    \end{column}
    \begin{column}{0.46\textwidth}
      \begin{takeaway}
        \small
        A clean diagnostic, usable on your own project:
        \vspace{0.3em}

        strong on multiple-choice and verifier benchmarks but weak with an
        LLM judge or on an arena $\Rightarrow$ the model knows the answer and
        cannot say it well $\Rightarrow$ this is the one case where preference
        optimisation is the right tool.
      \end{takeaway}
    \end{column}
  \end{columns}
  \sources{Raschka, \textit{From DeepSeek V3 to V3.2}, December 2025;
  \textit{Understanding the 4 Main Approaches to LLM Evaluation}.
  Model facts as of October 2026.}
  \note{The diagnostic in the box is genuinely useful and it is a nice
  hand-off to Jussi next week, because it only makes sense once you have four
  kinds of evaluation.

  The honest summary of this section: DPO made preference optimisation cheap
  enough that everyone does it, and then RLVR took over the half of the problem
  where you can check the answer. Preference methods now handle the parts that
  come down to taste.

  Clock check: should be around 17:00 here. Section 5 is the heart of the
  lecture.}
\end{frame}

% =====================================================================
\section{RL with verifiable rewards: GRPO}
% =====================================================================

\begin{frame}{The observation that produced reasoning models}
  \begin{block}{In one sentence}
    If you can \key{check} the answer automatically, you do not need to model
    human preference at all, which removes the part of post-training that
    refused to scale.
  \end{block}
  \vspace{0.5em}
  \begin{columns}[T]
    \begin{column}{0.52\textwidth}
      SFT needs written responses and RLHF needs preference labels. You buy
      both from people, at a fixed price per item.
      \vspace{0.4em}

      A verifier removes most per-item human labelling, but generation and
      checking still cost compute. Post-training therefore shifts much of the
      budget from annotation to computation.
    \end{column}
    \begin{column}{0.48\textwidth}
      \begin{block}{Names and dates}
        \small
        \textbf{RLVR} coined in \textbf{Tülu 3}, roughly two months before
        DeepSeek-R1.\\[0.4em]
        \textbf{DeepSeek-R1} and \textbf{Kimi K1.5} were posted to arXiv on
        the same day, \textbf{22 January 2025}.\\[0.4em]
        R1 became the famous one because of a single result: reasoning
        behaviour from \emph{pure} RL.
      \end{block}
    \end{column}
  \end{columns}
  \source{Raschka, \textit{Build Your Own Reasoning Model}, October 2026;
  \textit{The State of LLMs 2025}.}
  \note{Deliver the box slowly. This is the second of the three things to
  remember.

  The economics framing is what makes it click for most people: annotation
  budgets scale linearly with people, compute budgets scale with money, and
  the field had already proven it could convert money into capability.

  The simultaneous-arXiv anecdote is a good thirty seconds on how research
  credit actually works.}
\end{frame}

\begin{frame}{GRPO: one prompt, a group of rollouts, one update}
  \centering
  \begin{tikzpicture}
    \node[mod, minimum width=16mm] (p) at (0,0) {prompt $x$};
    \node[mod, minimum width=24mm] (r1) at (3.2, 1.45) {$y_1$\quad $r_1 = 1$};
    \node[mod, minimum width=24mm] (r2) at (3.2, 0.48) {$y_2$\quad $r_2 = 0$};
    \node[mod, minimum width=24mm] (r3) at (3.2,-0.48) {$y_3$\quad $r_3 = 1$};
    \node[mod, minimum width=24mm] (r4) at (3.2,-1.45) {$y_4$\quad $r_4 = 0$};
    \node[mod, minimum width=20mm] (g) at (6.4,0) {group stats\\$\bar{r}=0.5$};
    \node[modhi, minimum width=24mm] (u) at (9.3,0) {one update\\$A_i = (r_i-\bar{r})/s$};
    \draw[flow] (p) -- (r1);
    \draw[flow] (p) -- (r2);
    \draw[flow] (p) -- (r3);
    \draw[flow] (p) -- (r4);
    \draw[flow] (r1) -- (g);
    \draw[flow] (r2) -- (g);
    \draw[flow] (r3) -- (g);
    \draw[flow] (r4) -- (g);
    \draw[flowhi] (g) -- (u);
    \node[lbl] at (3.2, 2.15) {sample $G$ rollouts, verify each};
  \end{tikzpicture}

  \vspace{0.4em}
  \raggedright
  \small
  \textbf{1.} Sample $G$ complete responses for the \emph{same} prompt
  (``rollout'' is jargon for the full generated sequence).
  \textbf{2.} Verify each: 1 if correct, 0 if not.
  \textbf{3.} Score each response relative to the rest of \emph{its own group}.
  \textbf{4.} Clipped policy-gradient update.
  \source{Raschka, \textit{Build Your Own Reasoning Model}, October 2026,
  videos 6--7.}
  \note{Walk the diagram and emphasise the fan-out and collapse: $G$
  generations, one gradient step. Typical $G$ is 4 to 16; the cost of the
  method is almost entirely in the generation, not the backward pass.

  Make sure the word rollout is defined out loud. Half the room will not have
  seen it, and it appears in every RL paper they are about to read.}
\end{frame}

\begin{frame}{GRPO uses the group as its baseline}
  \[
    A_i \;=\; \frac{r_i \;-\; \text{mean reward of the group}}
                   {\text{standard deviation of the group}}
  \]
  \emph{In words:} the \key{advantage} says how much better an answer was than
  the other answers to the same prompt.
  \begin{itemize}
    \item PPO trains a critic to predict this baseline. GRPO measures it from
      the group, so there is no critic to hold or train.
    \item Example: rewards $1,0,1,0$ give advantages
      $+0.87,-0.87,+0.87,-0.87$.
    \item If every answer receives the same reward, every advantage is zero
      and that group teaches nothing.
  \end{itemize}
  \note{The last sentence of the box is the setup for the DAPO slide at the end
  of this section, so say it clearly now: a group where everything is right, or
  everything is wrong, produces no gradient.

  This is also the point at which to say that the advantage is assigned to the
  whole sequence. Every token in a correct rollout gets the same credit,
  including the irrelevant ones. Crude as that is, it works well enough.}
\end{frame}

% Optional mathematical and current-research detail.
\iffalse
\begin{frame}{The objective}
  \small
  \[
    \mathcal{J}_{\mathrm{GRPO}} = \mathbb{E}\!\left[
      \frac{1}{G}\sum_{i=1}^{G} \frac{1}{|y_i|} \sum_{t=1}^{|y_i|}
      \min\!\Big(\rho_{i,t} A_i,\;
      \mathrm{clip}(\rho_{i,t}, 1-\epsilon, 1+\epsilon)\, A_i\Big)
    \right]
    \;-\; \beta_{\mathrm{KL}}\,
    \mathbb{D}_{\mathrm{KL}}\!\left[\pi_\theta \,\|\, \pi_{\mathrm{ref}}\right]
  \]
  \normalsize
  \emph{In words:} average over the $G$ answers in the group and over the
  tokens of each one; nudge every token up or down in proportion to its
  answer's advantage, with a cap on how far a single batch may move the model;
  then subtract a penalty for drifting away from the frozen starting model.
  \vspace{0.3em}

  \begin{columns}[T]
    \begin{column}{0.5\textwidth}
      \small
      Compare with the PPO objective from twenty minutes ago. The clipped part
      is \emph{identical}, and only two things have changed: the critic is
      gone, and the advantage now comes from the group.
    \end{column}
    \begin{column}{0.5\textwidth}
      \small
      $\rho_{i,t}$ is again ``how much more likely is this token now than when
      we generated it''. Three models have roles here, not two: the one being
      trained, the one that \emph{generated} the answers (a few steps behind),
      and the frozen starting point.
      \vspace{0.3em}

      Both the KL term and the two divisions are contested --- last slide of
      this section.
    \end{column}
  \end{columns}
  \note{Do not derive this. Point at the three pieces they already know: the
  clipped ratio from PPO, the group advantage from the previous slide, and the
  KL leash from RLHF.

  The ratio in full, taken off the slide:
  $\rho_{i,t} = \pi_\theta(y_{i,t}\mid x, y_{i,<t}) /
  \pi_{\theta_{\mathrm{old}}}(y_{i,t}\mid x, y_{i,<t})$, which is the
  probability the current model assigns to token $t$ of answer $i$, divided by
  the probability the sampling model gave the same token in the same context. The
  $1/G$ averages over the group; the $1/|y_i|$ averages over the tokens of one
  answer, and it is this second division that Dr.\ GRPO removes.

  The old-versus-reference distinction trips people up every year, so say it
  twice. $\pi_{\theta_{\mathrm{old}}}$ keeps you near the model that produced
  the data; $\pi_{\mathrm{ref}}$ keeps you near the model you started from.}
\end{frame}
\fi

\begin{frame}{DeepSeek-R1-Zero: no SFT at all}
  \begin{columns}[T]
    \begin{column}{0.54\textwidth}
      \textbf{R1-Zero} applied RLVR \key{directly to the pretrained base
      model}. No supervised fine-tuning, no demonstrations, no preference data.
      \vspace{0.5em}

      The reward was two things added together, \textbf{is the answer right}
      and \textbf{is it laid out correctly}, plus a language-consistency term
      to stop the model alternating between languages mid-answer.
      \vspace{0.5em}

      \textbf{R1} itself is the practical model, trained in several stages from
      an SFT checkpoint of DeepSeek-V3. R1-Zero was the weaker of the two and
      by far the more important.
    \end{column}
    \begin{column}{0.46\textwidth}
      \begin{block}{What each reward does}
        \small
        \textbf{accuracy} --- is the final boxed answer right? One or zero,
        from a symbolic checker or a compiler.

        \textbf{format} --- is the reasoning inside the thinking tags and the
        answer where the parser expects it? This is the \emph{only} reason the
        tags exist.

        \textbf{language consistency} --- a patch for a real failure, where
        the model drifted between languages because nothing penalised it.
      \end{block}
    \end{column}
  \end{columns}
  \sources{Raschka, \textit{Build Your Own Reasoning Model}, October 2026;
  \textit{From DeepSeek V3 to V3.2}, December 2025. DeepSeek-R1 paper,
  January 2025.}
  \note{The reward written out, if you want it on the board:
  $R_{\mathrm{total}} = R_{\mathrm{accuracy}} + R_{\mathrm{format}}$, with the
  language-consistency term added as a third summand in the R1 (not R1-Zero)
  runs. All three are computed by code, not by a network, which is the whole
  point of the section.

  Contrast with Tülu 3 and Kimi K1.5, which both applied RL on top of an
  SFT model. RLVR was not R1's novelty. What R1 showed is that the supervised
  stage is not required for reasoning to appear.

  Worth saying: the thinking tags do nothing for reasoning ability. They are
  delimiters induced by a format reward so that the pipeline and the UI can
  separate the trace from the answer, and a different pair of strings would
  work just as well.}
\end{frame}

\begin{frame}{The counter-intuitive result}
  \begin{gotcha}[The reasoning trace is not directly supervised]
    In R1-Zero, the reward depends on the \textbf{final answer} and the
    \textbf{format}. The intermediate reasoning is not scored, not supervised,
    and not used to update the model. DeepSeek tried using it and reported that
    it did not help.
    \vspace{0.3em}

    Backtracking, self-checking, trying a second approach and abandoning a
    failed one all \textbf{emerge anyway}, from outcome rewards alone.
  \end{gotcha}
  \vspace{0.5em}
  \begin{columns}[T]
    \begin{column}{0.5\textwidth}
      \small
      \textbf{Why it works.} Traces that happen to contain a self-correction
      reach the right answer more often. Reinforcing the \emph{whole sequence}
      reinforces whatever made it succeed, and that includes the habit of
      checking.
    \end{column}
    \begin{column}{0.5\textwidth}
      \small
      \textbf{The part that should worry you.} Nothing in the objective asks
      the trace to be \emph{true}. It is optimised to precede correct answers,
      which is not the same thing. We return to this in two slides.
    \end{column}
  \end{columns}
  \source{Raschka, \textit{Build Your Own Reasoning Model}, October 2026,
  citing the DeepSeek-R1 paper.}
  \note{Gotcha three of three, and the best one. Ask the room to predict the
  answer before you reveal it: if you only reward the final answer, does the
  model learn to show its working? Most will say no.

  The right-hand column is the honest caveat and it sets up the reward-hacking
  slide, so do not skip it for the sake of a tidier story.}
\end{frame}

\begin{frame}[fragile]{How self-correction can emerge from outcome rewards}
\begin{lstlisting}
Q: Give me two numbers that sum to 10 and multiply to 21.

<think>
Let me try 5 and 5. The sum is 10. The product is 25.
Wait - that is not 21. Let me re-check what I need.
I need a + b = 10 and a * b = 21.
The factor pairs of 21 are (1, 21) and (3, 7).
1 + 21 = 22, no. 3 + 7 = 10, yes.
Check the product: 3 * 7 = 21. Both conditions hold.
</think>

3 and 7.
\end{lstlisting}
  \vspace{0.2em}
  \small
  The model starts down a wrong path, \key{notices}, says so, and recovers.
  Nobody wrote that behaviour into a dataset and nobody rewarded the
  recovery. It was reinforced because traces that contain it end in correct
  answers more often.
  \note{Read the trace out loud, including the ``wait''. It gets the reaction you
  want and it is the clearest piece of evidence for the previous slide.

  Then puncture it slightly. The word ``wait'' is not evidence of
  introspection. It is a token that, in this distribution, tends to come before
  a correction. Treat the behaviour as real and the anthropomorphic reading as
  unearned.

  Clock check: should be around 17:08 here.}
\end{frame}

\begin{frame}{Verifiers in practice}
  \footnotesize
  \centering
  \begin{tabular}{@{}L{2.6cm} L{4.4cm} L{4.2cm}@{}}
    \toprule
    \textbf{Domain} & \textbf{Checker} & \textbf{Where it breaks}\\
    \midrule
    mathematics & SymPy or WolframAlpha for symbolic equivalence &
      extraction and normalisation: is $0.5$ the same answer as
      $\tfrac{1}{2}$?\\ \addlinespace[0.25em]
    code & a compiler, unit tests, or a platform like LeetCode &
      tests that pass for the wrong reason; timeouts; sandbox escapes\\ \addlinespace[0.25em]
    computer use (2026) & success or failure signals from the environment &
      building the environments in the first place\\ \addlinespace[0.25em]
    structured output & a JSON-schema or grammar check &
      valid syntax, wrong content\\
    \bottomrule
  \end{tabular}

  \vspace{0.5em}
  \raggedright
  \small
  The pipeline is always the same three steps:
  \key{extract the final answer} $\rightarrow$ \key{normalise it} $\rightarrow$
  \key{check equivalence}. Most of the engineering is in the first two.
  The same verifier then does two jobs, scoring the model on your benchmark
  and supplying the reward during training.
  \source{Raschka, \textit{Build Your Own Reasoning Model}, October 2026;
  \textit{GPT-6 Astra, Looped Transformers, and Hidden Reasoning}.}
  \note{Expand row three. Building a computer-use environment is the whole
  cost of that row, and the scale is worth quoting: GPT-6 Astra reportedly ran
  tens of thousands of Mac Minis and Mac Studios as macOS environments for
  exactly this purpose.

  The dual-use point is the practical one for this course. If they build a
  decent answer extractor for Lab 2's baseline, they already have the reward
  function for the GRPO bonus. Say it explicitly, because it makes the bonus
  look achievable.

  The ``where it breaks'' column is where the real work is. Every group that
  tries GRPO on GSM8K will spend most of its time on answer extraction rather
  than on the RL.}
\end{frame}

% Optional research detail; verifier limitations remain on the preceding slide.
\iffalse
\begin{frame}[shrink=14]{Reward hacking}
  \begin{columns}[T]
    \begin{column}{0.54\textwidth}
      From the DeepSeekMath-V2 team, worth quoting exactly:
      \vspace{0.3em}

      \textit{``correct answers don't guarantee correct reasoning''}. A model
      can reach the right answer \emph{through flawed logic or fortunate
      errors}.
      \vspace{0.5em}

      And the problem that outcome rewards cannot touch at all: many
      mathematical tasks, theorem proving above all, require a rigorous
      derivation rather than a number. There is no final answer to check.
      \vspace{0.5em}

      \textbf{Models cannot grade themselves.} Asked to generate and then
      assess its own proof in one pass, the generator \emph{tends to claim
      correctness even when an external verifier easily finds the flaw}.
    \end{column}
    \begin{column}{0.46\textwidth}
      \begin{block}{Why R1 dropped process reward models}
        \small
        A PRM scores intermediate steps rather than outcomes, which is exactly
        what you would want here.
        \vspace{0.3em}

        DeepSeek's verdict: useful for reranking top-$N$ responses or guiding
        search, but its \emph{advantages are limited compared to the additional
        computational overhead} in large-scale RL.
        \vspace{0.3em}

        And it is hard to tell whether an intermediate reward is \emph{itself}
        correct --- so the PRM becomes the next thing to hack.
      \end{block}
    \end{column}
  \end{columns}
  \sources{Raschka, \textit{From DeepSeek V3 to V3.2}, December 2025;
  \textit{Understanding the 4 Main Approaches to LLM Evaluation}.}
  \note{The general law: whatever reward you can compute, the optimiser will
  find the cheapest way to maximise it. Nothing is broken when that happens.
  That is RL doing its job.

  The self-grading result is the one to connect forward to week 4's
  LLM-as-a-judge material. Once the judge is also the generator, you have
  stopped having a judge.}
\end{frame}

% The process-reward case study and optimisation details stay in this block.
\begin{frame}{How process rewards came back: DeepSeekMath-V2}
  \begin{columns}[T]
    \begin{column}{0.52\textwidth}
      Three models instead of one:
      \begin{enumerate}
        \item a \textbf{proof generator}
        \item a \textbf{proof verifier}, used as the reward model, trained
          against scores annotated by human mathematicians:
          \textbf{1} rigorous, \textbf{0.5} sound but with minor errors,
          \textbf{0} fatally flawed
        \item a \textbf{meta-verifier} that checks the verifier and stops it
          inventing problems
      \end{enumerate}
      \vspace{0.4em}
      Effect of the meta-verifier: the verifier's proof-analysis quality rose
      from \textbf{0.85} to \textbf{0.96}, at the same score-prediction
      accuracy.
    \end{column}
    \begin{column}{0.48\textwidth}
      \begin{block}{The GAN analogy}
        \small
        The verifier is the discriminator, the generator is the generator, and
        each one pushes the other along. Better proofs force the verifier to
        get sharper, and a sharper verifier raises the bar again.
      \end{block}
      \vspace{0.3em}
      \begin{block}{The twist}
        \small
        At \emph{inference} they use a \textbf{single} model for both
        generation and verification. The separate verifier was essential for
        \emph{training}; by the end, the generator has internalised the rubric.
      \end{block}
    \end{column}
  \end{columns}
  \source{Raschka, \textit{From DeepSeek V3 to V3.2}, December 2025.
  DeepSeekMath-V2, 27 November 2025.}
  \note{The arc of this slide is the arc of the section. Outcome rewards work
  until you reach the frontier, and the fix turns out to be a better reward
  rather than a better optimiser.

  The twist at the bottom is the elegant part and worth a beat. Train against
  a stronger critic for long enough and you end up with a model that carries
  the critic inside it.}
\end{frame}

\begin{frame}[t,shrink=9]{Five fixes that a GRPO run usually needs}
  \centering
  \small
  \begin{tabular}{@{}L{1.9cm} L{4.6cm} L{4.5cm}@{}}
    \toprule
    \textbf{Source} & \textbf{Trick} & \textbf{Problem it fixes}\\
    \midrule
    DAPO & throw away groups whose answers all scored the same &
      they have zero advantages and teach nothing\\ \addlinespace[0.25em]
    DAPO & average the loss over all tokens in the batch &
      otherwise a token in a long answer counts for less\\
      \addlinespace[0.25em]
    DAPO & allow a bigger step upwards than downwards (``clip-higher'') &
      an equal cap collapses the model onto a single answer\\
      \addlinespace[0.25em]
    Dr.\ GRPO & stop dividing by the group's standard deviation and by the
      answer length &
      those divisions favour \textbf{long wrong answers}\\
      \addlinespace[0.25em]
    Yao et al. & reweight each token by how far the two engines disagree &
      the two engines report different probabilities\\
    \bottomrule
  \end{tabular}

  \vspace{0.4em}
  \raggedright
  \small
  Olmo 3 adopts essentially all of these. DeepSeek V3.2 keeps the original
  divisions instead, and has training replay the sampling choices the rollout
  made.
  \source{Raschka, \textit{From DeepSeek V3 to V3.2}, December 2025.
  DAPO: Yu et al.\ 2025; Dr.\ GRPO: Liu et al.\ 2025.}
  \note{Raschka's own report is the line to quote: with these modifications in
  place, bad updates no longer corrupt his training runs and he no longer needs
  to reload checkpoints periodically.

  Details off the slide. Row 1 also generates replacement groups, so the batch
  size stays put. Row 2 means averaging over the whole batch rather than within
  each answer first. Row 4's two divisions also over-weight the easiest and
  hardest questions. Row 5 caps how large the correction is allowed to get, and
  the disagreement is about the same token. DeepSeek V3.2's replay means the
  two engines cannot disagree in the first place, and it also uses a separate
  leash strength per domain rather than one global KL coefficient.

  The precise terms, in case they read the papers. Row 1 is ``dynamic sampling
  of zero-variance groups''. Row 2 is the token-level policy-gradient loss.
  Row 3 is clip-higher: an upper clip bound $1+\epsilon_{\mathrm{high}}$ set
  above the lower $1-\epsilon_{\mathrm{low}}$, whose purpose is to stop the
  policy's entropy collapsing. Row 5 is truncated importance sampling. What
  DeepSeek V3.2 replays is the MoE expert routing and the top-p sampling mask
  from the rollout, so that the training-time action space matches the
  sampling-time one exactly.

  The last trick is the one with a sequel. Log-prob mismatch between the
  inference engine and the training engine is a determinism problem, and in
  week 3 we take it apart properly as batch invariance.}
\end{frame}
\fi

\begin{frame}{Which is why GRPO is a bonus on Thursday}
  \begin{columns}[T]
    \begin{column}{0.56\textwidth}
      \textbf{The cost structure.} Every step generates $G$ complete responses
      before a single gradient is computed. Generation dominates, so a GRPO
      step is one to two orders of magnitude more expensive than an SFT step.
      \vspace{0.5em}

      \textbf{The fragility.} Five known corrections, none of them default in
      every library, and failure looks like a loss curve that is fine while
      the outputs collapse.
      \vspace{0.5em}

      \textbf{The prerequisite.} A verifier you trust, which in practice means
      an answer extractor you have tested.
    \end{column}
    \begin{column}{0.44\textwidth}
      \begin{takeaway}
        \small
        \textbf{Bonus B}, for GSM8K groups especially: GRPO with an
        exact-match reward, about 100 steps, as a \texttt{gpumedium} batch job.
        \vspace{0.4em}

        One hundred steps will not make a reasoning model, but it should show
        you a reward curve that moves, and that is what you are after here.
        \vspace{0.4em}

        Do the main exercise first.
      \end{takeaway}
    \end{column}
  \end{columns}
  \note{Be direct: this is an optional, interesting, time-consuming exercise
  and is not required for the final project. Groups that attempt it should
  budget an evening, not an hour.

  Set the expectation about what success looks like. A reward curve that rises
  from 0.1 to 0.2 over a hundred steps is a successful lab result.

  Clock check: should be around 17:15 here. Ten minutes on LoRA.}
\end{frame}

% =====================================================================
\section{Parameter-efficient training}
% =====================================================================

\begin{frame}[shrink=6]{LoRA}
  Freeze the original weight matrix and never touch it again. Learn a small
  correction built from two thin matrices, and add that instead:
  \[
    W' \;=\; W \;+\; \frac{\alpha}{r}\, B A
  \]
  \emph{In words:} a 2560-by-2560 weight matrix holds 6.5 million numbers.
  Instead of retraining all of them, train a 2560-by-16 matrix and a
  16-by-2560 matrix, multiply them and add the result. That is 80 thousand
  numbers.

  \begin{columns}[T]
    \begin{column}{0.52\textwidth}
      \small
      $A$ starts random and $B$ starts at \textbf{zero}, so the correction is
      nothing at all at step 0 and training begins where the base model was.
      \vspace{0.4em}

      Trainable parameters for a 4B model at rank 16, on the attention and MLP
      matrices: roughly \textbf{20M}, or about \textbf{0.5\%}.
    \end{column}
    \begin{column}{0.48\textwidth}
      \begin{block}{Three knobs}
        \small
        \textbf{rank} --- how much the correction can express. 8 to 64 covers
        almost everything.

        \textbf{alpha} --- a scale factor. Set it to twice the rank.

        \textbf{target modules} --- which matrices get an adapter. Attention
        only is the cheap default.
      \end{block}
    \end{column}
  \end{columns}
  \source{Hu et al.\ 2021, \textit{LoRA: Low-Rank Adaptation of Large Language
  Models}.}
  \note{Three details the slide no longer has room for. A rank of 16 caps how
  much the correction can say, which is the point of the whole method. Alpha is
  divided by the rank, so changing the rank does not quietly change your
  effective learning rate, which is where the convention of twice the rank
  comes from. And adding the MLP matrices to the target modules usually pays
  for itself.

  Spell out the bet, which is no longer on the slide either. LoRA assumes the
  change a fine-tune needs is a simple one, even though the matrix it corrects
  is not. Fine-tuning selects a behaviour the model already has, and selecting
  does not require a rich correction, which is why a rank of 16 is enough so
  often.

  Shapes, in the paper's notation: $W \in \mathbb{R}^{d\times k}$ is
  frozen, $B \in \mathbb{R}^{d\times r}$ and $A \in \mathbb{R}^{r\times k}$ are
  trained, and $r \ll \min(d,k)$. The product $BA$ has the same shape as $W$
  but rank at most $r$, which is the whole constraint and also the name.

  The zero-initialisation of B is the detail worth pausing on. It is why
  you can attach a LoRA adapter to a production model and lose nothing before
  the first gradient step, and it is why LoRA training is so forgiving.

  Raschka lists LoRA alongside DPO as the example of academic work with a huge
  impact-to-budget ratio. Worth a sentence for a room of master's students.}
\end{frame}

\begin{frame}{The memory budget, worked, for a 4B model}
  \begin{columns}[T]
    \begin{column}{0.5\textwidth}
      \textbf{Full fine-tuning, mixed precision}
      \vspace{0.3em}
      \centering
      \small
      \begin{tabular}{@{}L{3.0cm} r@{}}
        \toprule
        weights, bf16 & 2 bytes/param\\
        gradients, bf16 & 2 bytes/param\\
        Adam momentum, fp32 & 4 bytes/param\\
        Adam variance, fp32 & 4 bytes/param\\
        fp32 master copy & 4 bytes/param\\
        \midrule
        \textbf{total} & \textbf{16 bytes/param}\\
        \bottomrule
      \end{tabular}
      \vspace{0.5em}

      \raggedright
      4 billion parameters at 16 bytes each is \textbf{64\,GB},
      \emph{before} a single activation.
    \end{column}
    \begin{column}{0.5\textwidth}
      \textbf{LoRA, bf16}
      \vspace{0.3em}
      \centering
      \small
      \begin{tabular}{@{}L{4.2cm} r@{}}
        \toprule
        frozen base weights, bf16 & 8\,GB\\
        adapter weights (20M) & 0.04\,GB\\
        adapter gradients & 0.04\,GB\\
        Adam state for 20M & 0.16\,GB\\
        \midrule
        \textbf{total state} & $\approx$ \textbf{8.3\,GB}\\
        \bottomrule
      \end{tabular}
      \vspace{0.5em}

      \raggedright
      Plus activations --- a few GB at sequence length 1024 with gradient
      checkpointing. Comfortably under \textbf{20\,GB}.
    \end{column}
  \end{columns}
  \vspace{0.5em}
  \centering
  One GH200 has \key{96\,GiB}, so full fine-tuning a 4B model would just about
  fit. LoRA fits five times over, and it trains faster and forgets less.
  \note{Make them do the 16-bytes-per-parameter arithmetic with you. It is the
  same accounting as Thursday of week 1, and it is the number that decides
  which experiments are possible on a given card.

  The honest footnote: with 96 GiB you could full fine-tune a 4B model. We
  choose LoRA for the other reasons, namely speed, a smaller footprint when
  you want several runs in an afternoon, and less forgetting.}
\end{frame}

% Optional tuning detail.
\iffalse
\begin{frame}[shrink=17]{Choosing rank, alpha and targets}
  \centering
  \begin{tabular}{@{}L{2.4cm} L{3.4cm} L{5.0cm}@{}}
    \toprule
    \textbf{Situation} & \textbf{Rank} & \textbf{Reasoning}\\
    \midrule
    format or style change, a few thousand examples & $r = 8$--$16$ &
      you are selecting an existing behaviour; capacity is not the bottleneck\\[0.35em]
    new task, tens of thousands of examples & $r = 32$--$64$ &
      more to learn, and enough data to avoid overfitting the adapter\\[0.35em]
    a genuinely new domain or language & higher, or full fine-tuning &
      if the behaviour is not latent in the base model, a low-rank nudge cannot
      summon it\\
    \bottomrule
  \end{tabular}

  \vspace{0.8em}
  \raggedright
  \small
  \textbf{Practical order:} fix $\alpha = 2r$ and tune the learning rate first.
  LoRA tolerates learning rates roughly ten times higher than full fine-tuning,
  so $1\times10^{-4}$ is a reasonable start where full fine-tuning would use
  $1\times10^{-5}$.
  \vspace{0.3em}

  \textbf{Target modules:} start with the query, key, value and output
  projections. Adding the gate, up and down projections roughly triples the
  adapter and usually helps on harder tasks.
  \note{The row that matters conceptually is the third one. LoRA works because
  fine-tuning is selection, so when the thing you want is not in there to be
  selected, the low-rank assumption is the thing that fails.

  Tell them the three hyperparameters they are allowed to tune in Lab 2 are
  learning rate, epochs and rank, and that everything else in the config stays
  as shipped so the comparison stays honest.}
\end{frame}
\fi

\begin{frame}{Merging, and serving adapters}
  \begin{columns}[T]
    \begin{column}{0.5\textwidth}
      \textbf{Merge.} $W \leftarrow W + \frac{\alpha}{r}BA$ produces an
      ordinary checkpoint with no adapter and \textbf{no inference overhead}.
      \vspace{0.4em}
      \begin{itemize}
        \item merge for a single deployment and for clean benchmarking
        \item merging is lossy only in the trivial sense that you can no longer
          take it off
        \item merge in the \emph{same} precision you trained in, or you will
          introduce a difference you did not intend
      \end{itemize}
    \end{column}
    \begin{column}{0.5\textwidth}
      \textbf{Keep it separate.} An adapter is tens of megabytes. Modern
      servers hot-swap them per request against one shared base model in
      memory.
      \vspace{0.4em}
      \begin{itemize}
        \item one base model, many tenants, one GPU
        \item a small amount of extra compute per token
        \item this is how fine-tuning gets offered as a service
      \end{itemize}
      \vspace{0.3em}
      \small
      In week 3 you will serve your own adapter this way with vLLM.
    \end{column}
  \end{columns}
  \vspace{0.4em}
  \raggedright
  \small
  \textbf{For Lab 2:} keep the adapter rather than a merged 8\,GB checkpoint.
  It is small and can be loaded again if you reuse it in later labs.
  \note{The commit-the-adapter instruction is operational and it matters: if
  twenty-five students each push a merged 8 GB model, the course project's disk
  quota disappears in an afternoon.

  The multi-tenant serving point is a nice preview of week 3 and explains to
  them how commercial fine-tuning APIs can be as cheap as they are.}
\end{frame}

% Optional boundary case.
\iffalse
\begin{frame}[shrink=5]{When LoRA genuinely underperforms}
  \begin{columns}[T]
    \begin{column}{0.54\textwidth}
      LoRA is close to full fine-tuning when you are \key{selecting} a
      behaviour. It falls behind when you are \key{adding} one:
      \vspace{0.4em}
      \begin{itemize}
        \item continued pretraining on a new domain or a new language
        \item very large instruction sets, where there is genuinely more to
          learn than a rank-16 update can hold
        \item tasks needing a change in the model's internal representation
          rather than in its output behaviour
      \end{itemize}
      \vspace{0.4em}
      The symptom is a training loss that plateaus well above where a full
      fine-tune would sit, and that does not improve when you lower the
      learning rate.
    \end{column}
    \begin{column}{0.46\textwidth}
      \begin{takeaway}
        \small
        LoRA's limitation and its benefit come from the same property.
        \vspace{0.4em}

        A rank-16 update cannot overwrite much, which is why it cannot teach a
        new language and also why it forgets so much less than a full
        fine-tune.
        \vspace{0.4em}

        If LoRA is not enough for your task, you have learned something about
        your task.
      \end{takeaway}
    \end{column}
  \end{columns}
  \note{The takeaway box is the sentence to leave them with, because it turns a
  disappointing lab result into a finding they can write up.

  If anyone asks about rank 128 or 256: at that point you are paying most of
  the memory of a full fine-tune for most of the constraint of LoRA, and you
  should probably just full fine-tune.

  Clock check: should be around 17:25 here.}
\end{frame}
\fi

% =====================================================================
\section{Prompt, retrieve, or fine-tune?}
% =====================================================================

\begin{frame}{The decision}
  \centering
  \begin{tabular}{@{}L{3.4cm} L{2.2cm} L{5.4cm}@{}}
    \toprule
    \textbf{You want to change} & \textbf{Reach for} & \textbf{Because}\\
    \midrule
    behaviour, format, style, tone, output schema, verbosity &
      \key{fine-tuning} &
      gradient descent reshapes the conditional distribution\\
      \addlinespace[0.6em]
    facts that are private, changing, or need a citation & \key{retrieval} &
      weights make a terrible database: no updates, no provenance, no access
      control\\
      \addlinespace[0.6em]
    nothing yet; you are still working out the task &
      \key{prompting} &
      no training cost, instant iteration, and it is your baseline anyway\\
    \bottomrule
  \end{tabular}

  \vspace{0.5em}
  \raggedright
  These are not really alternatives. Most production systems are
  \key{a fine-tuned model answering over retrieved context}, which is weeks 2
  and 5 of this course combined.
  \note{Give the one-line version first and then the table. Behaviour means
  training, knowledge means retrieval, and if you do not know yet, prompt.

  Push back on the framing in the question itself. ``Fine-tune or RAG'' is a
  false dichotomy that exists because the two have separate tutorials, and
  the project they are building will end up using both.}
\end{frame}

% Cost details are deferred to the inference and RAG lectures.
\iffalse
\begin{frame}[shrink=16]{What each one costs}
  \centering
  \small
  \begin{tabular}{@{}L{2.0cm} L{2.9cm} L{2.9cm} L{2.9cm}@{}}
    \toprule
    & \textbf{Prompting} & \textbf{Retrieval} & \textbf{Fine-tuning}\\
    \midrule
    up-front cost & none & build and host an index & one GPU run, about 20
      minutes for Lab 2 on a GH200\\ \addlinespace[0.2em]
    per-request cost & re-sent on every call; few-shot examples cost tokens
      forever & a retrieval hop, plus the retrieved tokens &
      \textbf{zero} once merged\\ \addlinespace[0.2em]
    latency & grows with prompt length; prefix caching helps &
      tens of milliseconds, plus a longer prefill & unchanged\\ \addlinespace[0.2em]
    iteration speed & seconds & minutes & hours\\ \addlinespace[0.2em]
    updating a fact & edit the prompt & re-index & retrain\\
    \bottomrule
  \end{tabular}

  \vspace{0.5em}
  \raggedright
  At low volume, prompt. Once the task is stable and the volume is high, that
  one GPU run is paid back quickly in shorter prompts.
  \note{The last line is the sentence an engineer actually uses. The shape of
  the argument is that prompting costs nothing to build and something on every
  request, while fine-tuning costs one GPU run and nothing afterwards. Make the
  volume argument concrete: a twelve-shot prompt that you can delete after
  fine-tuning saves those tokens on every single request, forever. The
  retrieved tokens in the second row land in the context window, which is why
  they cost you on every call.

  Prefix caching is the thing that weakens the argument slightly, and we take
  it apart properly in week 3. Flag it now so that nobody is surprised later.}
\end{frame}

\fi
\begin{frame}{What each one cannot fix}
  \begin{columns}[T]
    \begin{column}{0.33\textwidth}
      \begin{block}{Prompting cannot}
        \small
        create a capability the model does not have.
        \vspace{0.3em}

        It also cannot make a small model reliable at a long multi-step
        format. It will comply for six turns and then drift.
      \end{block}
    \end{column}
    \begin{column}{0.33\textwidth}
      \begin{block}{Retrieval cannot}
        \small
        fix behaviour. If the model will not follow your output schema,
        putting more text in its context will not help.
        \vspace{0.3em}

        It also cannot fix a model that ignores the context it was given.
      \end{block}
    \end{column}
    \begin{column}{0.33\textwidth}
      \begin{block}{Fine-tuning is a poor fit for}
        \small
        changing, attributable facts. It can teach factual associations, but
        narrow factual training can also \textbf{increase} hallucination
        \cit{Gekhman et al.\ 2024}.
        \vspace{0.3em}

        It teaches the model to answer confidently in that shape, which is not
        the same as knowing.
      \end{block}
    \end{column}
  \end{columns}
  \vspace{0.6em}
  \centering
  \key{Use retrieval when facts must be current, attributable or easy to
  update.}
  \vspace{0.3em}

  \raggedright
  \small
  Teams make this mistake every week, and it is an expensive one. Somebody
  fine-tunes on an internal FAQ, the model learns the \emph{form} of a
  confident internal answer, and it starts inventing policy numbers in exactly
  that form.
  \note{Spend real time on the right-hand column. It is the most common
  mistake in industry and the one they are most likely to make after this
  course.

  The practical risk is worth stating: maximum-likelihood training may teach
  the style of a confident answer more reliably than exact factual recall.
  Retrieval keeps the source available at inference time and can be updated
  without retraining.}
\end{frame}

\begin{frame}{The order of operations}
  \begin{enumerate}
    \item \textbf{Write a zero-shot prompt.} Score it on development data.
      This is your initial baseline.
      \vspace{0.3em}
    \item \textbf{Improve the prompt.} Few-shot examples, an explicit output
      format, a system prompt. Score again. Most teams stop improving here
      long before the prompt is actually good.
      \vspace{0.3em}
    \item \textbf{Ask what the remaining errors are.} Missing \emph{knowledge}
      $\rightarrow$ retrieval. Wrong \emph{behaviour} or \emph{format}
      $\rightarrow$ fine-tuning. Both $\rightarrow$ both.
      \vspace{0.3em}
    \item \textbf{Lock the protocol, then use the held-out test set once.}
      Compare the selected prompt and fine-tuned model, including behaviour
      you were not trying to change.
  \end{enumerate}
  \vspace{0.5em}
  \begin{takeaway}
    Steps 1 and 2 are not a formality. Beating a bad prompt tells you nothing,
    and beating a good one gives you a result worth writing down.
  \end{takeaway}
  \note{This is the methodological spine of Thursday, so say it as a procedure
  rather than as advice.

  The error-analysis step is the one students skip. Tell them to read thirty
  wrong outputs before deciding what to do next. Thirty is enough to see the
  pattern and short enough to do in twenty minutes.}
\end{frame}

\begin{frame}{The Lab 2 question}
  \centering
  \vspace{0.6em}
  {\Large\key{For your task, does fine-tuning beat your best prompt?}}
  \vspace{1.0em}

  \raggedright
  Note everything that question requires of you:
  \begin{itemize}
    \item \textbf{your task} --- use the task you explored in Lab 1 or choose
      another suitable task, with its own held-out set
    \item \textbf{your best prompt} --- not your first prompt. Zero-shot and
      few-shot, on \emph{both} the base and the instruct model
    \item \textbf{beat} --- a number on data neither model was trained on,
      with identical decoding settings on both sides
  \end{itemize}
  \vspace{0.6em}
  There is one more, and it is the part most people leave out: \key{what did it
  cost?} That is why the table has four numbers rather than three.
  \vspace{0.5em}

  \small
  Forward pointers: \textbf{week 5} is the retrieval branch of the decision
  slide, \textbf{week 6} wraps a tool-using agent around it, and \textbf{week
  4} is how you find out whether any of these numbers mean anything.
  \note{Pause on this slide. It is the single sentence the whole of Thursday
  answers, and it is also a section of their final report.

  ``Not your first prompt'' deserves emphasis. A weak baseline is the easiest
  way to produce an impressive-looking and worthless result, and the rubric
  rewards the honest comparison.

  Clock check: should be around 17:35 here. Ten minutes of logistics left.}
\end{frame}

% =====================================================================
\section{Lab 2 briefing}
% =====================================================================

\begin{frame}[shrink=4]{Material for the final report}
  Keep one table with \key{four numbers}, one loss curve and one paragraph.
  \vspace{0.5em}
  \centering
  \begin{tabular}{@{}L{3.4cm} C{2.2cm} L{5.0cm}@{}}
    \toprule
    \textbf{Setting} & \textbf{Held-out score} & \textbf{Notes}\\
    \midrule
    zero-shot, instruct & & your floor\\
    best prompt & & which model, which prompt, how many shots\\
    LoRA fine-tuned & & rank, epochs, learning rate\\
    forgetting check & & 50 general-knowledge questions, before \emph{and} after\\
    \bottomrule
  \end{tabular}

  \vspace{0.8em}
  \raggedright
  \small
  Plus: the training loss curve, and one paragraph answering
  \emph{was it worth it?}, in GPU-minutes and in points.
  \vspace{0.3em}

  A negative result is fine. ``Fine-tuning lost to a six-shot prompt, and here
  is why we think so'' scores just as well, provided the numbers are real.
  \vspace{0.2em}

  \texttt{lab2-example/README.md} contains the runnable baseline, LoRA,
  evaluation and plotting workflow.
  \note{Put the empty table on screen and explain that it provides material
  for the final report. Showing the shape of the result at the briefing is worth
  more than any amount of description.

  Say the negative-result line explicitly. Every year a group panics in the
  last twenty minutes because the fine-tune did not win, and starts tuning
  hyperparameters against the held-out set to fix it.}
\end{frame}

\begin{frame}{The session, and what to do before it}
  \begin{columns}[T]
    \begin{column}{0.5\textwidth}
      \textbf{Thursday, 16:15--17:45}
      \vspace{0.3em}
      \small
      \begin{tabular}{@{}l L{4.0cm}@{}}
        \toprule
        0--10 & set-up, GPU session, model loaded\\
        10--35 & baselines: zero-shot and few-shot, base and instruct\\
        35--70 & LoRA SFT\\
        70--90 & measure, forgetting check, save results\\
        \bottomrule
      \end{tabular}
      \vspace{0.5em}

      \raggedright
      Thirty-five minutes for the fine-tune is not generous. Arrive with your
      data already formatted and you should get through it.
    \end{column}
    \begin{column}{0.5\textwidth}
      \begin{gotcha}[No homework compute today or tomorrow]
        Roihu is down for scheduled maintenance on \textbf{3--4 November}.
        Nothing queues, nothing runs, and the web interface is unavailable.
        \vspace{0.3em}

        It returns on \textbf{Wednesday evening}. Do the parts that need no GPU
        now: read your model's chat template, format your \texttt{messages}
        file, write the answer extractor.
      \end{gotcha}
    \end{column}
  \end{columns}
  \note{Say the maintenance window clearly and tell them it is not an excuse to
  arrive unprepared, because the preparation that matters is data formatting
  and that needs no GPU at all.

  Three things to do tonight, on a laptop: print the chat template, convert
  500 training examples to the conversational format, and write the function
  that extracts an answer from a generated string.}
\end{frame}

\begin{frame}{Baselines, and three ways to ruin one}
  \begin{columns}[T]
    \begin{column}{0.5\textwidth}
      \textbf{The protocol}
      \begin{itemize}
        \item zero-shot \emph{and} few-shot
        \item base \emph{and} instruct
        \item fixed decoding settings. Greedy removes sampling variation,
          although exact GPU repeatability is not guaranteed
        \item the held-out set is \textbf{never} touched during development
      \end{itemize}
      \vspace{0.3em}
      Choose the best of those four on dev. Report which one it was, then score
      that locked choice on the held-out test set.
    \end{column}
    \begin{column}{0.5\textwidth}
      \textbf{What goes wrong}
      \vspace{0.3em}
      \small
      \textbf{1. Answer extraction.} The model was right and your regular
      expression disagreed. Read twenty raw generations before you trust any
      score.\\[0.4em]
      \textbf{2. Chat template on a base model.} It has none. Prompt it as a
      document, few-shot, with a consistent separator.\\[0.4em]
      \textbf{3. Contamination.} Public benchmark, public model. Go and
      investigate a suspiciously high zero-shot score before you believe it.
    \end{column}
  \end{columns}
  \vspace{0.4em}
  \raggedright
  \small
  All three move the \emph{baseline}, which is the number everything else gets
  compared against, so an error here propagates through every remaining lab.
  \note{These three are worth more of the briefing than the fine-tuning itself.
  A broken fine-tune is obvious within minutes, whereas a broken baseline can
  sit there unnoticed for weeks.

  Give the concrete instruction: print twenty raw generations next to their
  extracted answers and read them. It takes five minutes and it catches the
  first failure mode outright.}
\end{frame}

% The maintained code is in lab2-example; avoid a dense code slide here.
\iffalse
\begin{frame}[fragile,shrink=17]{The fine-tune, concretely}
\begin{lstlisting}[language=Python]
from trl import SFTConfig, SFTTrainer
from peft import LoraConfig

peft_config = LoraConfig(
    r=16, lora_alpha=32, lora_dropout=0.05, task_type="CAUSAL_LM",
    target_modules=["q_proj","k_proj","v_proj","o_proj",
                    "gate_proj","up_proj","down_proj"])

args = SFTConfig(
    output_dir="out/lora-sft",
    assistant_only_loss=True,      # mask the prompt - check this worked
    bf16=True,                     # GH200 supports bf16 natively
    per_device_train_batch_size=4,
    gradient_accumulation_steps=4, # effective batch 16
    gradient_checkpointing=True,
    learning_rate=1e-4, num_train_epochs=2,
    lr_scheduler_type="cosine", warmup_ratio=0.03,
    max_length=1024, logging_steps=10)

trainer = SFTTrainer(model=model_id, args=args,
                     train_dataset=ds, peft_config=peft_config)
trainer.train()
\end{lstlisting}
  \vspace{0.2em}
  \small
  Dataset format: a \texttt{messages} column of role/content dicts. TRL applies
  \textbf{the model's own chat template}. You tune \textbf{three} things:
  learning rate, epochs, rank.
  \note{Walk the config, not the API. The lines that matter are
  assistant-only loss, bf16, and the effective batch size of sixteen from four
  times four.

  Tell them to verify the mask once by printing a batch's labels. The flag
  silently does nothing if the chat template does not mark the assistant span,
  and that failure costs an afternoon.

  Two epochs is a starting point, not a recommendation. With 1 to 5k examples,
  three epochs usually overfits visibly.}
\end{frame}

\fi
\begin{frame}{Measuring, optional extensions, and files to keep}
  \begin{columns}[T]
    \begin{column}{0.5\textwidth}
      \textbf{The forgetting check}
      \vspace{0.2em}
      \small
      Fifty general-knowledge multiple-choice questions, nothing to do with
      your task. Score the base model, score yours, and report the delta,
      using the same prompt and the same decoding both times.
      \vspace{0.6em}

      \normalsize\textbf{Bonuses}
      \vspace{0.2em}
      \small
      \textbf{A --- DPO} on a preference set. Start beta at 0.1, and remember
      that the reference model is just your adapter switched off.
      \vspace{0.2em}

      \textbf{B --- GRPO} with an exact-match reward, about 100 steps, as a
      \texttt{gpumedium} batch job. GSM8K groups especially.
    \end{column}
    \begin{column}{0.5\textwidth}
      \begin{block}{Keep for the final submission}
        \small
        the results table \textbullet\ the loss curve \textbullet\ the
        \textbf{adapter}, not a merged checkpoint \textbullet\ the notebook and
        the Slurm script \textbullet\ your prompt file
      \end{block}
      \vspace{0.2em}
      \begin{block}{Roihu reminders}
        \small
        No compute \textbf{3--4 November} --- maintenance.

        \texttt{gputest} before \texttt{gpumedium}.

        Close idle GPU Jupyter sessions; one forgotten session holds a whole
        GH200 for twelve hours.

        Build environments on the ARM login nodes.
      \end{block}
    \end{column}
  \end{columns}
  \note{The file list is also useful in week 3. The serving lab can load the
  adapter produced here, but students may instead use another compatible
  model or adapter.

  Repeat the maintenance dates one final time, and repeat the idle-session
  rule. Those two sentences prevent most of the support traffic in week 2.}
\end{frame}

% =====================================================================
% A title-only summary frame adds no content.
\iffalse
\begin{frame}[standout]
  Three things to remember
  \note{Deliberate pause. These three sentences are repeated in the week 7
  wrap-up, so slow down and let people write them.}
\end{frame}
\fi

\begin{frame}{Three things to remember}
  \begin{enumerate}
    \item \key{A chat is a string, and the chat template is the whole
      interface.} Roles, special tokens and the generation prompt are a format
      the model was trained to continue.
      \vspace{0.6em}
    \item \key{If you can verify the answer you do not need a reward model.}
      That observation turned post-training from an annotation budget into a
      compute budget, and it is what produced reasoning models.
      \vspace{0.6em}
    \item \key{Always measure what fine-tuning broke, not only what it
      improved.} GLM-5 needed a whole extra distillation stage to repair its
      own RL stages. Your fifty-question check is the same idea, for one
      afternoon's work.
  \end{enumerate}
  \vfill
  \small
  Across all three methods, post-training usually reshapes and strengthens
  behaviour with much less data than pretraining. It can teach task-specific
  information, but it is not a reliable store for changing facts.
  \note{Say all three slowly, then the closing claim. If the room remembers one
  sentence from today it should be the last one on this slide.

  Point back at the opening demo: the base model had everything it needed to
  answer that question. What it lacked was a reason to believe it had been
  asked.}
\end{frame}

\begin{frame}[shrink=3]{Thursday, and the next two weeks}
  \begin{columns}[T]
    \begin{column}{0.5\textwidth}
      \textbf{Thursday 5 November --- Lab 2}\\
      Does fine-tuning beat your best prompt? Bring formatted data and a
      working answer extractor.
      \vspace{0.8em}

      \textbf{Week 3 --- Aarne Talman, inference}\\
      Decoding, KV-cache memory, serving, quantisation and the cost of running
      your model.
    \end{column}
    \begin{column}{0.5\textwidth}
      \begin{block}{Why inference comes next}
        \small
        Today you changed a model. Next week you put a model behind a server
        and find out what it costs to run. Reusing today's adapter is useful,
        but not required.
        \vspace{0.4em}

        You will compare the original and FP8 versions with the same held-out
        metric used today, so speed and quality appear in one table.
      \end{block}
    \end{column}
  \end{columns}
  \vspace{0.4em}
  \raggedright
  \small
  \textbf{Week 4 --- Jussi Karlgren, evaluation:} metrics and their limits,
  contamination, LLM-as-a-judge and evaluation across languages.
  \note{Explain the new order clearly: week 2 studies behaviour change, week 3
  studies serving, and week 4 examines whether the evaluation itself is
  trustworthy. Reusing one model makes the comparison easier to follow, but
  students may change models when their task requires it.}
\end{frame}

\appendix

\begin{frame}[shrink=5]{Reading}
  \begin{columns}[T]
    \begin{column}{0.5\textwidth}
      \textbf{The pipeline}
      \begin{itemize}\small\setlength{\itemsep}{0pt}
        \item Ouyang et al.\ 2022, InstructGPT
        \item Lambert et al.\ 2024, \textit{Tülu 3}
        \item Zhou et al.\ 2023, \textit{LIMA}
      \end{itemize}
      \vspace{0.3em}
      \textbf{Preference optimisation}
      \begin{itemize}\small\setlength{\itemsep}{0pt}
        \item Schulman et al.\ 2017, PPO
        \item Rafailov et al.\ 2023, \textit{DPO}
        \item Ethayarajh et al.\ 2024, \textit{KTO}
        \item Hong et al.\ 2024, \textit{ORPO}
        \item Meng et al.\ 2024, \textit{SimPO}
      \end{itemize}
      \vspace{0.3em}
      \textbf{Parameter-efficient training}
      \begin{itemize}\small\setlength{\itemsep}{0pt}
        \item Hu et al.\ 2021, \textit{LoRA}
        \item Gekhman et al.\ 2024
      \end{itemize}
    \end{column}
    \begin{column}{0.5\textwidth}
      \textbf{RLVR and reasoning}
      \begin{itemize}\small\setlength{\itemsep}{0pt}
        \item DeepSeek-AI 2025, \textit{DeepSeek-R1}
        \item Yu et al.\ 2025, \textit{DAPO}
        \item Liu et al.\ 2025, \textit{Dr.\ GRPO}
        \item DeepSeekMath-V2, November 2025
      \end{itemize}
      \vspace{0.3em}
      \textbf{Current practice}, Raschka
      \begin{itemize}\small\setlength{\itemsep}{0pt}
        \item \textit{Build Your Own Reasoning Model} (Oct 2026)
        \item \textit{Controlling Reasoning Effort in LLMs} (Jul 2026)
        \item \textit{From DeepSeek V3 to V3.2} (Dec 2025)
        \item \textit{The State of LLMs 2025}
      \end{itemize}
      \vspace{0.3em}
      \textbf{Libraries}
      \begin{itemize}\small\setlength{\itemsep}{0pt}
        \item TRL docs: SFT, DPO, KTO, GRPO trainers
      \end{itemize}
    \end{column}
  \end{columns}
  \note{If they read two things, make it the DPO paper for the derivation and
  the R1 paper for the RLVR story, including its section on unsuccessful
  attempts, which is more instructive than the successful ones.

  Worth saying out loud what each one is for, since the slide only gives
  names. InstructGPT is the original RLHF recipe. Tülu 3 coins RLVR and is the
  most complete open post-training report we have, so it is the one to point
  at for anybody who wants to reproduce a full pipeline rather than read about
  one. LIMA is the quality-over-quantity result for SFT. DeepSeek-R1 carries
  GRPO, R1-Zero and the aha moments. DeepSeekMath-V2 is where the verifier and
  meta-verifier come from. Gekhman et al.\ is the paper on fine-tuning with
  new knowledge and the hallucination it causes.}
\end{frame}

\end{document}
